% Backing appendices for the recap deck (modules/recap/slides): one chapter
% per claim behind the deck's design (bilaga B-E), and bilaga F, the
% flipped-classroom review adapted from the course-evaluation report.  Written by the recap-claims agent,
% 2026-10-10, to be \input after the deck's own method chapter (bilaga A).
% Bibliography: ltnotes-repetition.bib (merge into the deck's ltnotes.bib).
% Hit-list tables: litteratursokning/recap-*-full.tex.
% Labels are prefixed app:recap-, tab:recap- and sec:recap- for book-wide
% uniqueness.  The chapters refer to the deck only in general terms ("den
% här föreläsningen"); the orchestrator may add \cref's to the deck's
% sections where marked "% POINT:".

% Let the generated longtables start on the page where they are input.
\setcounter{LTchunksize}{4}

\chapter{Lär man sig mer av att plocka fram svaret själv än av att läsa
  det igen?}
\label{app:recap-retrieval}
\chapterprecis{Författaren har ännu inte granskat resultaten i den här
  bilagan i sin helhet.}

Ja, med förbehåll: att plocka fram svaret ur minnet ger bättre minne på
sikt än att läsa om det, också i verkliga kurser och mer med
återkoppling, men vinsten gäller mest det som frågan gällde, den är liten
och osäker i matematik, och den är inte prövad med kontrollerad design i
programmeringskurser.

Varje avsnitt i den här föreläsningen börjar med en fråga som du svarar
på innan svaret visas, den första i \cref{ovn:recap-bankomat}.
Skälet är det som kallas testeffekten (\foreignlanguage{english}{testing
effect}): att plocka fram något ur minnet, en hämtningsövning
(\foreignlanguage{english}{retrieval practice}), ger bättre minne på sikt
än att läsa samma sak en gång till\autocite{Rowland2014Testing}.  Det är
ett empiriskt påstående, och det har tre delar som måste prövas var för
sig: att effekten finns, att den håller i en verklig kurs med korta
frågor utan betyg --- inte bara i laboratoriet --- och var den tar slut.
Frågan den här bilagan besvarar är: \emph{lär sig studenter i en
verklig kurs mer av att svara på frågor om det de redan mött än av att
läsa om det, och under vilka villkor gäller det?}

\section{Metod}
\label{sec:recap-retrieval-metod}

Arbetsgången är den i föreläsningen \emph{Algoritmiskt tänkande},
bilaga~A.  Frågan söktes ämnesdrivet och tvåsidigt den 10 oktober 2026
med verktyget \texttt{scholar}: tre stödfrågor --- om effekten i
klassrummet, i programmering och närliggande ämnen, och i
översiktsstudier --- och två motfrågor som letade efter gränser, uteblivna
effekter och arbeten om överföring till nya frågor och om återkoppling.
Motfrågorna använder samma tre namn på begreppet som stödfrågorna.  Varje
fråga ställdes till OpenAlex, IEEE Xplore, Web of Science och Scopus.
Semantic Scholar avvisade sin nyckel, och DBLP svarade den dagen med en
kontrollsida mot robotar i stället för träffar på varje fråga, så de två
databaserna saknas.  \Cref{tab:recap-retrieval-queries} visar frågorna
och antalet träffar.  En språkmodell sorterade träffarna som stödjande,
kvalificerande, angränsande eller felträffar; varje källa som citeras
nedan är läst och vald av en människa, i fulltext där en fri kopia fanns
och annars i sammanfattningen, vilket anges i källans provenansanteckning
och i \cref{sec:recap-retrieval-diskussion}.

% Author's note (verification and reproduction):
%   scholar session : recap-retrieval (S1-S3, M1-M2 on openalex, ieee,
%                     wos, scopus, -l 40; DBLP word lists D1-D4 sent but
%                     answered by DBLP's anti-bot page, 0 rows)
%   query strings   : claims-work/claims-queries.tsv in the recap-claims
%                     scratch directory; counts in claims-counts.tsv
%   searched        : 2026-10-10; providers checked with `scholar
%                     providers check` (s2: key rejected, HTTP 403)
%   classification  : scholar llm context + llm classify --no-examples,
%                     model openai-codex/gpt-5.6-sol; cited rows decided
%                     by hand with theme tags
%   provenance      : ltnotes-repetition.bib, keys Yang2021Testing
%                     Agarwal2021Retrieval Rowland2014Testing
%                     PanRickard2018Transfer vanGogSweller2015
%                     KarpickeAue2015 Rawson2015Complex Trumbo2021Testing
%                     Pan2019Jargon McDaniel2011Quiz DonoghueHattie2021
%                     Alzaid2017BiteSize Wojtowicz2025SQL
%                     Glaser2025LectureHall Schwerter2026Gateway
%                     Murray2025Mathematics; vanAlten2019 HewLo2018
%                     (from bilaga F)
%   note            : Rowland2014Testing and DonoghueHattie2021 surfaced
%                     in the recap-spacing session (bilaga C), not here.
\begin{table}[htbp]
  \centering
  \caption{Sökfrågor och antal träffar per databas, 10 oktober 2026,
    verktyget \texttt{scholar}.  S = stödfråga, M = motfråga.  R står
    för (\enquote{retrieval practice} \textsc{or} \enquote{testing
    effect} \textsc{or} \enquote{test-enhanced learning}).  Databaser:
    OpenAlex (OA; inte \foreignlanguage{english}{open access}), IEEE
    Xplore (IEEE), Web of Science (WoS), Scopus.  Högst 40 träffar
    hämtades per databas och fråga, så 40 betyder \enquote{minst 40}.
    Frågan visas i den booleska form som skickades till OpenAlex och IEEE
    Xplore; Web of Science fick samma sträng som \texttt{TS=(\dots)} och
    Scopus som \texttt{TITLE-ABS-KEY(\dots)}.  DBLP och Semantic Scholar
    svarade inte (se texten).}
  \label{tab:recap-retrieval-queries}
  \footnotesize
  \begin{tabular}{@{}lp{0.52\linewidth}rrrr@{}}
    \toprule
    & Nyckelord & OA & IEEE & WoS & Scopus \\
    \midrule
    S1 & R \textsc{and} (classroom \textsc{or} course \textsc{or} lecture
      \textsc{or} \enquote{real-world})
      & 40 & 40 & 40 & 40 \\
    S2 & R \textsc{and} (programming \textsc{or} \enquote{computer
      science} \textsc{or} STEM \textsc{or} engineering)
      & 40 & 40 & 40 & 40 \\
    S3 & R \textsc{and} (\enquote{meta-analysis} \textsc{or}
      \enquote{systematic review})
      & 39 & 0 & 40 & 40 \\
    M1 & R \textsc{and} (complexity \textsc{or} \enquote{complex
      materials} \textsc{or} \enquote{boundary conditions} \textsc{or}
      \enquote{no benefit} \textsc{or} \enquote{null effect} \textsc{or}
      limitations)
      & 40 & 9 & 40 & 40 \\
    M2 & R \textsc{and} (transfer \textsc{or} feedback \textsc{or}
      \enquote{low-stakes} \textsc{or} quizzing) \textsc{and} (moderator
      \textsc{or} moderators \textsc{or} \enquote{does not} \textsc{or}
      fails)
      & 40 & 4 & 40 & 40 \\
    \bottomrule
  \end{tabular}
\end{table}

\section{Resultat}
\label{sec:recap-retrieval-resultat}

\paragraph{Finns effekten?}
Ja, och den är en av de bäst belagda i inlärningsforskningen.  En
metaanalys av försök där ett test jämförs med att läsa om samma
material, mest i laboratoriet, fann en tillförlitlig fördel för testet,
\enquote{g = 0.50 (CI [0.42, 0.58])}, med stor spridning mellan
studierna \parencite{Rowland2014Testing}.  En senare metaanalys av tio
studieteknikers effekter fann att just spridd övning och övningstest
hade de största effekterna av alla\autocite{DonoghueHattie2021}.

\paragraph{Håller den i en verklig kurs?}
Ja.  Den största genomgången som bara tar med klassrumsstudier --- 222
studier med 48\,478 elever och studenter --- fann att frågor i
undervisningen höjer resultaten med ungefär en halv standardavvikelse,
\enquote{g = 0.499}, och med \enquote{g = 0.330} jämfört med att läsa
om\autocite{Yang2021Testing}.  Samma genomgång fann att återkoppling
efter frågan ökar vinsten (\enquote{g = 0.537} med och \enquote{g =
0.374} utan), att det inte spelar någon roll om frågan besvaras med
penna, på nätet eller med ett svarssystem i salen, och att frågor med
eller utan betydelse för betyget gav samma vinst \parencite{Yang2021Testing}.
En systematisk genomgång med en snävare definition --- bara försök i
riktiga klasser där varje student svarade själv --- fann medelstora
eller stora vinster i 57 procent av 49 effektstorlekar, och negativa i
bara tre\autocite{Agarwal2021Retrieval}.  Ett försök i riktiga klasser
som flyttade frågorna mellan före föreläsningen, efter den och som
repetition före provet fann den största vinsten för repetitionsfrågorna,
och vinsten fanns kvar på terminens och årets
slutprov\autocite{McDaniel2011Quiz}.
Också i flippade klassrum höjer quiz på förberedelserna effekten,
\(g = 0{,}19\)\autocite[9]{vanAlten2019}, och effekten är större när
passet inleds med quiz\autocite{HewLo2018}%
% Bilaga F is in the teacher's edition only.
\ifbool{ltnote}{; den litteraturen redovisas i
\cref{app:recap-flippat}}{}.
I en kurs i neuroanatomi gav korta frågor med svarsapparat i salen
bättre provresultat upp till tolv veckor senare\autocite{Pan2019Jargon}.

\paragraph{Gäller det programmering?}
Det vet vi inte, eftersom ingen kontrollerad studie av testeffekten i
en programmeringskurs fanns bland träffarna.  Klassrumsmetaanalysen
innehåller bara två datavetenskapliga kurser \parencite{Yang2021Testing}.
De två studier som ligger närmast visar samband, inte effekter: de
studenter som använde ett verktyg med korta programmeringsfrågor mer
klarade kursen bättre\autocite{Alzaid2017BiteSize}, och de som gjorde
fler frivilliga övningsprov i en kurs i databasfrågor (SQL) fick högre
betyg\autocite{Wojtowicz2025SQL} --- men i båda valde studenterna själva
hur mycket de övade.  Det närmaste ämnet med kontrollerade studier är
matematik, och där är bilden splittrad.  I en grundkurs i matematik med
omkring 1\,800 studenter, där övningsuppgifterna lottades, gav
hämtningsövningar 4--5 procentenheter bättre resultat på delproven,
mest på likadana uppgifter och minst på uppgifter som krävde överföring,
och på slutprovet bara i en av tre årskullar\autocite{Schwerter2026Gateway}.
En metaanalys av matematikinlärning fann testeffekten där liten och
osäker, \enquote{g = 0.18. However, the 95\% confidence interval crossed
zero}\autocite{Murray2025Mathematics}.

\paragraph{Var tar effekten slut?}
Motsökningen gav tre förbehåll.  Det första gäller komplext material:
van Gog och Sweller hävdar att effekten avtar och kan försvinna
\enquote{when the complexity of learning material is very
high}\autocite{vanGogSweller2015}.  Det påståendet är omstritt --- två
svar i samma temanummer menar att data tvärtom visar \enquote{a small
positive effect of retrieval practice with worked
examples}\autocite{KarpickeAue2015} och att slutsatsen inte är
underbyggd\autocite{Rawson2015Complex} --- men alla tre är eniga om att
komplext material är för lite studerat.  Programmering är just sådant
material.  Det andra gäller överföring: det man lärt sig genom att
svara på en fråga hjälper även på andra frågor, \enquote{d = 0.40},
men minst för material som inte själv testades och för problem med
lösta exempel, och utan gynnsamma villkor ofta inte
alls\autocite{PanRickard2018Transfer}.
Kursen i neuroanatomi visar samma sak i liten skala: frågorna hjälpte på
provfrågor om samma begrepp, men inte konsekvent på frågor som krävde
förståelse \parencite{Pan2019Jargon}, och i föreläsningar i psykologi gav
övningsfrågor efter föreläsningen en effekt för det som övades men ingen
alls för det som togs upp på samma föreläsningar utan att
övas\autocite{Glaser2025LectureHall}.  Det tredje gäller vad
resultaten mäter: de flesta studierna mäter fakta och ytlig kunskap, och
författarna till översikten över tio tekniker manar själva till
försiktighet när resultaten förs över till djupare
förståelse \parencite{DonoghueHattie2021}.  En översikt från laboratoriet
till klassrummet sammanfattar läget: effekten är robust, men om den
gynnar mer avancerade studenter är inte undersökt\autocite{Trumbo2021Testing}.

\section{Diskussion och begränsningar}
\label{sec:recap-retrieval-diskussion}

Beläggen för att effekten finns och håller i klassrummet är starka och
kommer från oberoende håll: två metaanalyser och en systematisk
genomgång med olika urval.  Svagheten ligger i det som skiljer den här
kursen från de studerade: programmering är komplext material, frågorna
i föreläsningen är ofta förutsägelser om ett programs beteende snarare
än fakta, och ingen av källorna gäller nybörjare i Python.

Flera källor är lästa bara i sammanfattningen eftersom ingen fri kopia
fanns: metaanalysen om överföring \parencite{PanRickard2018Transfer}, de
två svaren i debatten om komplext material
\parencite{KarpickeAue2015,Rawson2015Complex},
översikten från laboratoriet till klassrummet \parencite{Trumbo2021Testing},
försöket med frågornas placering \parencite{McDaniel2011Quiz}, studien i
föreläsningar i psykologi \parencite{Glaser2025LectureHall}, studien och
metaanalysen i matematik \parencite{Schwerter2026Gateway,Murray2025Mathematics}
och de två studierna nära programmering
\parencite{Alzaid2017BiteSize,Wojtowicz2025SQL}.  Klassrumsmetaanalysen
är läst i en version utgiven före tryckningen, utan sidnummer.  Två av
sex databaser svarade inte.

\section{Slutsats}
\label{sec:recap-retrieval-slutsats}

Ja, med förbehåll: att svara på frågor om det man mött ger bättre minne
på sikt än att läsa om det, också i verkliga kurser och med frågor utan
betydelse för betyget, med en effekt kring en halv standardavvikelse,
större med återkoppling och oberoende av om svaret ges med telefon eller
penna.  Förbehållen är att effekten är mindre säker för komplext
material --- i matematik är den liten och osäker --- att den mest gäller
just det som frågades om och knappt sprider sig till resten, och att den
aldrig är prövad med kontrollerad design i en programmeringskurs.
Föreläsningen
kan därför säga att frågorna hjälper minnet och att svaret ska följas av
återkoppling, men inte att de ger förståelse av sig själva --- därför
följs varje fråga av exempel och genomgång.

\section{Fortsatt arbete}
\label{sec:recap-retrieval-fortsatt}

Sökningen saknar Semantic Scholar och DBLP; särskilt DBLP indexerar de
datavetenskapliga konferenser där en programmeringsstudie skulle ha
publicerats, så sökningen S2 bör göras om där.  Tio källor är lästa
bara i sammanfattningen (se \cref{sec:recap-retrieval-diskussion}), och
metaanalysen av Adesope med flera från 2017 dök bara upp som en
rättelse\footnote{Rättelsen har DOI 10.3102/0034654317700823.} och är
inte läst.

Vad fältet saknar är ett kontrollerat försök i en programmeringskurs: en
del av innehållet repeteras med frågor och en del med en vanlig
genomgång, och båda delarna prövas på provet med uppgifter som kräver
att studenten skriver eller förklarar kod, inte bara känner igen den.
Visar ett sådant försök att frågorna hjälper också där, kan föreläsningen
motivera dem med programmeringens egna data; visar det att effekten
försvinner för att skriva kod, bör frågorna ses som en hjälp att minnas
begreppen och kompletteras med skrivövningar.

\needspace{10\baselineskip}
\section{Träffar som rör påståendet}

\Cref{tab:sok-recap-retrieval} listar de träffar som rör påståendet.

% Author's note: generated with `scholar sessions export recap-retrieval
% -f table --bearing-only --lang sv --label sok-recap-retrieval --track
% "..." --theme ... -o litteratursokning/recap-retrieval-full`.
% --- scholar LaTeX fragment --------------------------------------------
% Generated by: scholar sessions export -f table
% Session: recap-retrieval
% \input this file from the body of a document whose
% preamble loads:
%   \usepackage{longtable}
%   \usepackage{booktabs}
%   \usepackage{array}
% ----------------------------------------------------------------------

\begingroup\footnotesize
\begin{longtable}{@{}%
  >{\raggedright\arraybackslash}p{0.44\textwidth}c%
  >{\raggedright\arraybackslash}p{0.13\textwidth}%
  >{\raggedright\arraybackslash}p{0.26\textwidth}@{}}
\caption{Träffar som rör påståendet, hämtningsövning och testeffekten (14 citerade, 66 som stöder påståendet och 126 som kvalificerar eller motsäger det, av 409 unika träffar; de 150 angränsande och 53 felträffarna listas inte). Skälet till varje inkludering står i sista kolumnen; för maskinklassade rader anges modellens konfidens. Databaser: OpenAlex (OA; inte open access), IEEE Xplore (IEEE), Web of Science (WoS), Scopus.}\label{tab:sok-recap-retrieval}\\
\toprule
Titel & År & Leverantör & Skäl \\
\midrule
\endfirsthead
\multicolumn{4}{@{}l}{\emph{\tablename~\thetable{} (forts.)}}\\
\toprule
Titel & År & Leverantör & Skäl \\
\midrule
\endhead
\midrule \multicolumn{4}{r@{}}{\emph{forts.\ på nästa sida}}\\
\endfoot
\bottomrule
\endlastfoot
A Meta-analytic Review of the Effectiveness of Spacing and Retrieval Practice for Mathematics Learning. & 2025 & OA, Scopus & \emph{citerad}: programmering och matematik \\
Is the testing effect ready to be put to work? Evidence from the laboratory to the classroom. & 2021 & OA & \emph{citerad}: gränser: komplexitet och överföring \\
Not New, but Nearly Forgotten: the Testing Effect Decreases or even Disappears as the Complexity of Learning Materials Increases & 2015 & OA & \emph{citerad}: gränser: komplexitet och överföring \\
Online and Clicker Quizzing on Jargon Terms Enhances Definition-Focused but Not Conceptually Focused Biology Exam Performance & 2019 & WoS & \emph{citerad}: gränser: komplexitet och överföring \\
Retrieval Practice Consistently Benefits Student Learning: a Systematic Review of Applied Research in Schools and Classrooms & 2021 & Scopus & \emph{citerad}: i klassrummet \\
Retrieval Practice in Higher Education: Causality and Content Transfer Effects in a Gateway Math Course & 2026 & WoS & \emph{citerad}: programmering och matematik \\
Test-enhanced learning in a middle school science classroom: The effects of quiz frequency and placement. & 2011 & OA & \emph{citerad}: i klassrummet \\
Testing (Quizzing) Boosts Classroom Learning: A Systematic an Meta-Analytic Review & 2021 & WoS & \emph{citerad}: i klassrummet \\
The effects of bite-size distributed practices for programming novices & 2017 & IEEE & \emph{citerad}: programmering och matematik \\
The Status of the Testing Effect for Complex Materials: Still a Winner & 2015 & OA & \emph{citerad}: gränser: komplexitet och överföring \\
The Testing Effect in the Lecture Hall: Does it Transfer to Content Studied but Not Practiced? & 2025 & Scopus & \emph{citerad}: gränser: komplexitet och överföring \\
The Testing Effect Is Alive and Well with Complex Materials & 2015 & OA & \emph{citerad}: gränser: komplexitet och överföring \\
Train Hard, Score High: The Impact of Practice Tests on SQL Learning & 2025 & Scopus & \emph{citerad}: programmering och matematik \\
Transfer of test-enhanced learning: Meta-analytic review and synthesis & 2018 & Scopus & \emph{citerad}: gränser: komplexitet och överföring \\
A dual memory theory of the testing effect & 2017 & OA & stöder påståendet (0.96) \\
A Large Language Model-Based System for Socratic Inquiry: Fostering Deep Learning and Memory Consolidation & 2025 & Scopus & stöder påståendet (0.75) \\
A Multimodal Human--AI Instructional Framework for Productive Vocabulary Development: A Classroom Evaluation of a Coordinated LLM--ASR System & 2026 & Scopus & stöder påståendet (0.81) \\
A Systematic Review of the Testing Effect in Learning & 2013 & OA & stöder påståendet (1.00) \\
Achieving Testing Effects in an Authentic College Classroom & 2021 & OA & stöder påståendet (0.96) \\
Adaptive spaced-repetition training, working memory, and long-term recall: A quasi-experimental study in university students & 2026 & Scopus & stöder påståendet (0.98) \\
AI Teaches Surgical Diagnostic Reasoning to Medical Students: Evidence from an Experiment Using a Fully Automated, Low-Cost Feedback System & 2025 & WoS & stöder påståendet (0.82) \\
Association between ScholarRx Qmax question bank engagement and NBME CBSE and USMLE Step 1 performance & 2026 & WoS & stöder påståendet (0.70) \\
Brain network differences in second language learning depend on individual competencies & 2026 & WoS & stöder påståendet (0.84) \\
Chatbots for vocabulary learning: A mixed-methods study with A1-level learners in a French as a foreign language course & 2026 & Scopus & stöder påståendet (0.76) \\
Comparing the testing effect under blocked and mixed practice: The mnemonic benefits of retrieval practice are not affected by practice format & 2016 & OA & stöder påståendet (0.99) \\
Corrigendum to: Rethinking the use of tests: A meta-analysis of practice testing (Review of Educational Research, (2017), 10.3102/0034654316689306) & 2017 & Scopus & stöder påståendet (0.68) \\
Covert Retrieval Practice Benefits Retention As Much As Overt Retrieval Practice & 2011 & OA & stöder påståendet (0.98) \\
Covert retrieval yields a forward testing effect across levels of successive list similarity & 2023 & Scopus & stöder påståendet (0.96) \\
Does Individual Performance Feedback Increase the Use of Retrieval Practice? & 2021 & OA & stöder påståendet (0.93) \\
Ecological Validity of the Testing Effect & 2016 & OA & stöder påståendet (0.99) \\
Effects of reinforcement on test-enhanced learning in a large, diverse introductory college psychology course. & 2016 & OA & stöder påståendet (1.00) \\
Enhancing Electrical Science Learning Within A Novel Practice Testing Learning Framework & 2018 & IEEE & stöder påståendet (0.87) \\
Enhancing Final Exam Performance Through Retrieval Practice: Evidence From a Diverse College Sample & 2026 & Scopus & stöder påståendet (0.99) \\
Enhancing visuospatial learning: The benefit of retrieval practice & 2010 & OA & stöder påståendet (0.99) \\
Extension of the dual-memory model of test-enhanced learning to distributions and individual differences & 2020 & OA & stöder påståendet (0.95) \\
Far transfer of retrieval-practice benefits: rule-based learning as the underlying mechanism & 2024 & WoS & stöder påståendet (0.99) \\
Gamified digital mediation and the development of American English phonetic competence in secondary EFL learners & 2026 & WoS & stöder påståendet (0.76) \\
Group discussions and test-enhanced learning: individual learning outcomes and personality characteristics & 2017 & Scopus & stöder påståendet (0.98) \\
Improving Mathematical Problem Solving in an Introductory Engineering Course With the Testing Effect & 2019 & OA & stöder påståendet (0.91) \\
Improving retention by placing retrieval practice at the end of class: a naturalistic study & 2020 & OA & stöder påståendet (0.99) \\
Improving students' summative knowledge of introductory chemistry through the forward testing effect: examining the role of retrieval practice quizzing & 2020 & OA & stöder påståendet (0.95) \\
Incentives for retrieval practice and exam performance of college students & 2026 & WoS & stöder påståendet (0.94) \\
LAVA-Era Learning Method Infrastructure: An AI-Supported Socratic and Spaced Practice System & 2026 & IEEE & stöder påståendet (0.71) \\
Learning Verbs in Sentences: Children With Developmental Language Disorder and the Role of Retrieval Practice & 2024 & Scopus & stöder påståendet (0.99) \\
Low-stakes Quizzing Enhances Learning in Two Undergraduate Plant Identification Courses & 2026 & Scopus & stöder påståendet (0.96) \\
Metacognition of the testing effect: Guiding learners to predict the benefits of retrieval & 2012 & OA & stöder påståendet (0.90) \\
Metacognitive judgments can potentiate new learning: The role of covert retrieval & 2022 & OA & stöder påståendet (0.93) \\
Multiple-choice testing as a desirable difficulty in the classroom. & 2014 & OA & stöder påståendet (0.98) \\
Neural Signatures of Test-Potentiated Learning in Parietal Cortex & 2013 & OA & stöder påståendet (0.98) \\
Optimizing Japanese Kanji learning through Anki-assisted spaced repetition: the mediating role of cognitive engagement and visual-verbal integration efficiency & 2026 & Scopus & stöder påståendet (0.87) \\
Practicing Retrieval Facilitates Learning & 2020 & OA & stöder påståendet (1.00) \\
Quizzing and retention in the high school science class & 2013 & OA & stöder påståendet (0.96) \\
Regaining access to marginal knowledge in a classroom setting & 2020 & WoS & stöder påståendet (0.98) \\
Reinforcing General Chemistry Learning with Structured Daily Recaps: A Pilot Study & 2026 & Scopus & stöder påståendet (0.82) \\
Retrieval practice - a tool to be able to retain higher mathematics even 3 months after the exam & 2024 & WoS & stöder påståendet (0.98) \\
Retrieval practice and spacing in an engineering mathematics classroom: Do the effects add up? & 2017 & IEEE & stöder påståendet (0.99) \\
Retrieval Practice for Effective Teaching of an Engineering Course & 2021 & IEEE & stöder påståendet (0.87) \\
Retrieval Practice Improves Exam Performance as a Function of Review Question Number and Format & 2025 & WoS & stöder påståendet (0.99) \\
Retrieval practice improves memory in patients with schizophrenia: new perspectives for cognitive remediation & 2019 & OA & stöder påståendet (0.91) \\
Retrieval Practice in Classroom Settings: A Review of Applied Research & 2019 & OA & stöder påståendet (1.00) \\
Retrieval Practice Promotes Semantic Integration of Real-World Memories Without Compromising Episodic Integrity & 2025 & OA, Scopus & stöder påståendet (0.92) \\
Retrieval practice with feedback facilitates the acquisition of English spelling rules and etymological knowledge in EFL learners & 2026 & Scopus & stöder påståendet (0.99) \\
Sleep reduces the testing effect---But not after corrective feedback and prolonged retention interval. & 2018 & OA & stöder påståendet (0.99) \\
Spaced Retrieval Practice in Two Fundamental Engineering Courses: Calculus and Physics & 2022 & IEEE & stöder påståendet (0.78) \\
STEMing evidence-based study strategies: a systematic review of literature & 2026 & Scopus & stöder påståendet (0.96) \\
Taking the testing effect beyond the college freshman: Benefits for lifelong learning. & 2013 & OA & stöder påståendet (0.98) \\
Test-enhanced learning in the classroom: Long-term improvements from quizzing. & 2011 & OA & stöder påståendet (0.99) \\
Testing beyond words: Using tests to enhance visuospatial map learning & 2007 & OA & stöder påståendet (1.00) \\
Testing during study insulates against the buildup of proactive interference. & 2008 & OA & stöder påståendet (0.95) \\
Testing effect in L2 discourse comprehension: importance of retrieval-based learning & 2024 & WoS & stöder påståendet (0.99) \\
Testing Enhances the Transfer of Learning & 2012 & OA & stöder påståendet (0.99) \\
The ``Testing Effect'' in Undergraduate Endodontic Education: The Impact of Regular Low Stakes Testing on Learning Outcomes and Student Engagement & 2025 & OA & stöder påståendet (0.98) \\
The critical role of retrieval practice in long-term retention & 2010 & OA & stöder påståendet (1.00) \\
The Effect of Question Placement on Learning from Textbook Chapters & 2018 & Scopus & stöder påståendet (1.00) \\
The effects of tests on learning and forgetting & 2008 & OA & stöder påståendet (0.99) \\
The role of retrieval mode and retrieval orientation in retrieval practice: insights from comparing recognition memory testing formats and restudying & 2016 & OA & stöder påståendet (0.99) \\
The Testing Effect in the Psychology Classroom: A Meta-Analytic Perspective & 2017 & OA & stöder påståendet (1.00) \\
Using Retrieval Practice on Learning with Children: Systematic Literature Review & 2022 & Scopus & stöder påståendet (0.98) \\
What types of learning are enhanced by a cued recall test? & 2006 & OA & stöder påståendet (0.99) \\
WIP: Analyzing Students' Practices Behaviors in an Introductory Computer Science Course and Monitoring Their Practices Behaviors in a Subsequent Class & 2024 & Scopus & stöder påståendet (0.76) \\
``Yes we can?'' -- To what extent do can-judgments enhance self-regulated learning from text? & 2026 & Scopus & kvalificerar eller motsäger påståendet (0.82) \\
A (Preliminary) Recipe for Obtaining a Testing Effect in Preschool Children: Two Critical Ingredients & 2018 & OA & kvalificerar eller motsäger påståendet (0.99) \\
A Classroom Study on the Role of Prior Knowledge and Retrieval Tool in the Testing Effect & 2020 & OA & kvalificerar eller motsäger påståendet (0.98) \\
A Meta-Analytic Review of the Benefit of Spacing out Retrieval Practice Episodes on Retention & 2021 & Scopus & kvalificerar eller motsäger påståendet (0.95) \\
A role for familiarity in supporting the testing effect over time & 2019 & OA & kvalificerar eller motsäger påståendet (0.96) \\
A Scoping Review of Retrieval Practice (Test-Enhanced Learning) in Nursing Education & 2023 & Scopus & kvalificerar eller motsäger påståendet (0.97) \\
A Study of the Testing Effect in an Engineering Classroom & 2020 & OA & kvalificerar eller motsäger påståendet (0.87) \\
Accuracy Matters for the Benefits of Sleep After Retrieval Practice & 2020 & OA & kvalificerar eller motsäger påståendet (0.95) \\
Active recall strategies associated with academic achievement in young adults: A systematic review & 2024 & Scopus & kvalificerar eller motsäger påståendet (0.94) \\
Adaptive retrieval practice with multiplechoice questions in the university classroom & 2020 & OA & kvalificerar eller motsäger påståendet (1.00) \\
Adding the keyword mnemonic to retrieval practice: A potent combination for foreign language vocabulary learning? & 2019 & OA & kvalificerar eller motsäger påståendet (0.99) \\
Anki Use and Academic Performance in Medical Education: A Systematic Review of Evidence and Learning Theory & 2026 & Scopus & kvalificerar eller motsäger påståendet (0.88) \\
Application of Retesting as a Learning Tool in an Online Science Course & 2024 & Scopus & kvalificerar eller motsäger påståendet (0.92) \\
Assessing Boundary Conditions of the Testing Effect: On the Relative Efficacy of Covert vs. Overt Retrieval & 2017 & OA & kvalificerar eller motsäger påståendet (0.99) \\
Boundary Conditions of Retrieval-Enhanced Learning: Evidence from Contextual Learning and Retention of Second Language Multi-Word Expressions & 2026 & Scopus & kvalificerar eller motsäger påståendet (0.99) \\
Cloze tests as retrieval practice activities: evaluating their integration with audience response systems in K-12 schools & 2025 & Scopus & kvalificerar eller motsäger påståendet (1.00) \\
Cognitive perspectives on maintaining physicians' medical expertise: IV. Best practices and open questions in using testing to enhance learning and retention & 2023 & Scopus & kvalificerar eller motsäger påståendet (0.99) \\
Collaborative and Individual Practice Testing Lead to Different Patterns of Learning Across Two Variations of the True/False Test Format & 2026 & WoS & kvalificerar eller motsäger påståendet (0.99) \\
Comparing Form-focused Retrieval and Meaning-inferencing Approaches to Phrasal Verb Learning by Korean EFL Learners & 2026 & Scopus & kvalificerar eller motsäger påståendet (0.97) \\
Diminishing-cues retrieval practice: A memory-enhancing technique that works when regular testing doesn't & 2017 & OA & kvalificerar eller motsäger påståendet (1.00) \\
Do differences in topic knowledge matter? An experimental investigation into topic knowledge as a possible moderator of the testing effect & 2025 & Scopus & kvalificerar eller motsäger påståendet (0.99) \\
Does Covert Retrieval Benefit Adolescents' Learning in 8th Grade Science Classes? & 2025 & Scopus & kvalificerar eller motsäger påståendet (0.99) \\
Does pre-testing promote better retention than post-testing? & 2019 & OA & kvalificerar eller motsäger påståendet (0.99) \\
Does response mode affect amount recalled or the magnitude of the testing effect? & 2012 & OA & kvalificerar eller motsäger påståendet (0.97) \\
Does Retrieval Demand Moderate the Effectiveness of Covert Retrieval Practice? Comparing Covert and Overt Retrieval Practice for Learning Key Terms and Definitions & 2025 & Scopus & kvalificerar eller motsäger påståendet (1.00) \\
Does retrieval practice depend on semantic cues? Assessing the fuzzy trace account of the testing effect & 2017 & OA & kvalificerar eller motsäger påståendet (0.99) \\
Does test-enhanced learning transfer for triple associates? & 2015 & OA & kvalificerar eller motsäger påståendet (0.99) \\
Effects of case revisitation during deliberate reflection on medical students' diagnostic accuracy: a randomized experiment & 2026 & WoS & kvalificerar eller motsäger påståendet (0.96) \\
Effects of random selection tests on second language vocabulary learning: a comparison with cumulative tests & 2025 & Scopus & kvalificerar eller motsäger påståendet (0.91) \\
Effects of sleep and retrieval practice on verbal paired-associate learning across 12 and 24 hours & 2023 & OA & kvalificerar eller motsäger påståendet (0.98) \\
Elaboration and the Testing Effect in Cued Recall & 2010 & OA & kvalificerar eller motsäger påståendet (0.90) \\
Encouraging Knowledge Transfer in Food Science and Nutrition Education: Suggestions from Cognitive Research & 2019 & Scopus & kvalificerar eller motsäger påståendet (0.86) \\
Enhancing Clinical Reasoning in Surgical Education: A Randomized Crossover Trial on AI-Generated Individual Feedback for Key Feature Question-Based Online Assignments & 2026 & WoS & kvalificerar eller motsäger påståendet (0.94) \\
Enhancing learning and retarding forgetting: Choices and consequences & 2007 & OA & kvalificerar eller motsäger påståendet (0.94) \\
Errors May Not Cue Recall of Corrective Feedback: Evidence Against the Mediation Hypothesis of the Testing Effect & 2021 & WoS & kvalificerar eller motsäger påståendet (0.95) \\
Evidence and Theory for why the Best Example-Problem Ratio To Optimize Learning Gain Depends on Knowledge Content & 2025 & WoS & kvalificerar eller motsäger påståendet (0.99) \\
Examining the benefits of spaced retrieval practice on transfer and the student experience in the classroom & 2026 & Scopus & kvalificerar eller motsäger påståendet (1.00) \\
Examining the effect size and duration of retrieval-induced facilitation & 2023 & Scopus & kvalificerar eller motsäger påståendet (0.91) \\
Examining the Testing Effect in University Teaching: Retrievability and Question Format Matter & 2018 & OA & kvalificerar eller motsäger påståendet (1.00) \\
Example-based learning: should learners receive closed-book or open-book self-explanation prompts? & 2020 & WoS & kvalificerar eller motsäger påståendet (0.94) \\
Exploring the Mechanisms that Underlie the Benefits of Retrieval Practice in Younger and Older Adults & 2021 & OA & kvalificerar eller motsäger påståendet (0.90) \\
Formats and implementations of exercises for collocation learning: Learning outcomes and students' beliefs & 2025 & WoS & kvalificerar eller motsäger påståendet (0.94) \\
Gamified feedback in adaptive retrieval practice: Points and progress-bars enhance motivation but not learning & 2026 & Scopus & kvalificerar eller motsäger påståendet (0.94) \\
Generalizing test-enhanced learning from the laboratory to the classroom & 2007 & OA & kvalificerar eller motsäger påståendet (1.00) \\
Happy or redundant? Retrieval practice enhances lasting learning in non-interactive teaching for low-quality explanations & 2027 & Scopus & kvalificerar eller motsäger påståendet (0.98) \\
Impact of different practice testing methods on learning outcomes & 2024 & WoS & kvalificerar eller motsäger påståendet (0.96) \\
In search of transfer following cued recall practice: The case of process-based biology concepts & 2019 & Scopus & kvalificerar eller motsäger påståendet (1.00) \\
Interaction of repetition and retention interval influences the transfer effect after answer feedback for episodic memory & 2025 & Scopus & kvalificerar eller motsäger påståendet (0.98) \\
Interpolated Retrieval of Relevant Material, Not Irrelevant Material, Enhances New Learning of a Video Lecture In-Person and Online & 2025 & Scopus & kvalificerar eller motsäger påståendet (0.99) \\
Investigating the testing effect: Retrieval as a characteristic of effective study strategies & 2018 & OA & kvalificerar eller motsäger påståendet (0.98) \\
Is Covert Retrieval an Effective Learning Strategy? Is It as Effective as Overt Retrieval? Answers from a Meta-Analytic Review & 2025 & Scopus & kvalificerar eller motsäger påståendet (1.00) \\
Learning by writing explanations: Is explaining to a fictitious student more effective than self-explaining? & 2021 & WoS & kvalificerar eller motsäger påståendet (0.93) \\
Learning techniques to improve memory in children: a systematic review & 2025 & Scopus & kvalificerar eller motsäger påståendet (0.98) \\
Levels of Retrieval and the Testing Effect & 2021 & WoS & kvalificerar eller motsäger påståendet (0.99) \\
Long-term inference and memory following retrieval practice & 2020 & Scopus & kvalificerar eller motsäger påståendet (0.99) \\
Lower constraint testing enhances the testing effect for some contextual details but not others & 2024 & OA & kvalificerar eller motsäger påståendet (0.94) \\
Mechanisms behind the testing effect: an empirical investigation of retrieval practice in meaningful learning & 2015 & OA & kvalificerar eller motsäger påståendet (0.97) \\
Mixed effects of retrieval practice on the acquisition of procedural knowledge in verb spelling education & 2026 & WoS & kvalificerar eller motsäger påståendet (0.98) \\
Mnemonic benefits of retrieval practice at short retention intervals & 2014 & OA & kvalificerar eller motsäger påståendet (0.97) \\
Motivation Brought to the Test: Mastery Goal Orientation and External Rewards for Successful Retrieval Practice Modulate the Testing Effect for Complex Study Contents & 2023 & OA & kvalificerar eller motsäger påståendet (0.99) \\
Motivation brought to the test: Successful retrieval practice is modulated by mastery goal orientation and external rewards & 2023 & OA & kvalificerar eller motsäger påståendet (1.00) \\
Multiple Practice Success Scaffolds Long-Term Test-Enhanced Learning in Preschoolers & 2025 & Scopus & kvalificerar eller motsäger påståendet (0.99) \\
Optimizing schedules of retrieval practice for durable and efficient learning: How much is enough? & 2011 & OA & kvalificerar eller motsäger påståendet (0.97) \\
PLAT 20(1) 2021: Enhancing Student Learning in Research and Educational Practice: The Power of Retrieval Practice and Feedback & 2021 & Scopus & kvalificerar eller motsäger påståendet (0.92) \\
Providing corrective feedback during retrieval practice does not increase retrieval-induced forgetting & 2013 & Scopus & kvalificerar eller motsäger påståendet (0.98) \\
Research on \textquotedbl{}the Testing Effect\textquotedbl{} Routinely Conflates Direct and Forward Testing Effects: A Meta-Analysis of Testing Effects With Free Recall & 2026 & WoS & kvalificerar eller motsäger påståendet (1.00) \\
Response format, not semantic activation, influences the failed retrieval effect & 2019 & Scopus & kvalificerar eller motsäger påståendet (0.91) \\
Retrieval Effects on Memory: Overview of Current Basic Research & 2021 & WoS & kvalificerar eller motsäger påståendet (0.96) \\
Retrieval enhanced learning: Does it matter what topic is to be learned? & 2026 & WoS & kvalificerar eller motsäger påståendet (0.99) \\
Retrieval practice enhances near but not far transfer of spatial memory & 2020 & Scopus & kvalificerar eller motsäger påståendet (1.00) \\
Retrieval Practice in Memory- and Language-Impaired Populations: A Systematic Review & 2020 & OA & kvalificerar eller motsäger påståendet (0.96) \\
Retrieval Practice in Stepwise Worked Examples Improves Learning & 2025 & WoS & kvalificerar eller motsäger påståendet (0.99) \\
Retrieval practice may not benefit mathematical word-problem solving & 2023 & Scopus & kvalificerar eller motsäger påståendet (1.00) \\
Retrieval Practice Versus Elaborative Encoding: A Systematic and Meta-analytic Review & 2025 & Scopus & kvalificerar eller motsäger påståendet (1.00) \\
Retrieval Practice: Beneficial for All Students or Moderated by Individual Differences? & 2020 & OA & kvalificerar eller motsäger påståendet (0.99) \\
Retrieval-induced versus restudy-induced forgetting in serial order memory & 2026 & Scopus & kvalificerar eller motsäger påståendet (0.91) \\
Reversing the testing effect by feedback is a matter of performance criterion at practice & 2020 & OA & kvalificerar eller motsäger påståendet (0.99) \\
Reversing the testing effect by feedback: Behavioral and electrophysiological evidence & 2016 & OA & kvalificerar eller motsäger påståendet (0.99) \\
Scaling Retrieval Practice with LLM: Improving Multiple Choice Question (MCQ) Quality through Knowledge Graphs & 2026 & Scopus & kvalificerar eller motsäger påståendet (0.96) \\
Selective retrieval practice can benefit long-term retention of untested associations: Evidence from a word triplet paradigm & 2024 & OA & kvalificerar eller motsäger påståendet (0.96) \\
Sleep can reduce the testing effect: It enhances recall of restudied items but can leave recall of retrieved items unaffected. & 2014 & OA & kvalificerar eller motsäger påståendet (1.00) \\
Spaced Retrieval Practice Improves Engineering Student Performance in Physics & 2023 & IEEE & kvalificerar eller motsäger påståendet (0.97) \\
Systematic review of distributed practice and retrieval practice in health professions education & 2024 & Scopus & kvalificerar eller motsäger påståendet (0.99) \\
Techniques for scaffolding retrieval practice: The costs and benefits of adaptive versus diminishing cues & 2019 & OA & kvalificerar eller motsäger påståendet (0.99) \\
Test-Enhanced Learning & 2023 & OA & kvalificerar eller motsäger påståendet (0.97) \\
Test-Enhanced Learning and Incentives in Biology Education & 2020 & OA & kvalificerar eller motsäger påståendet (0.99) \\
Test-enhanced learning in health professions education: A systematic review: BEME Guide No. 48 & 2018 & OA & kvalificerar eller motsäger påståendet (0.99) \\
Test-Enhanced Learning: The Potential for Testing to Promote Greater Learning in Undergraduate Science Courses & 2015 & OA & kvalificerar eller motsäger påståendet (0.99) \\
Testenhanced learning in medical education & 2008 & OA & kvalificerar eller motsäger påståendet (1.00) \\
Testing and transfer: Retrieval practice effects across test formats in English vocabulary learning in school & 2021 & OA & kvalificerar eller motsäger påståendet (0.99) \\
Testing improves performance as well as assesses learning: A review of the testing effect with implications for models of learning. & 2022 & OA & kvalificerar eller motsäger påståendet (0.97) \\
Testing Memories of Personally Experienced Events: The Testing Effect Seems Not to Persist in Autobiographical Memory & 2018 & OA & kvalificerar eller motsäger påståendet (0.97) \\
Testing the testing effect on prolific: when retrieval practice fails to boost learning & 2026 & Scopus & kvalificerar eller motsäger påståendet (0.99) \\
Testing the testing effect with advanced materials while accounting for individual differences & 2026 & Scopus & kvalificerar eller motsäger påståendet (1.00) \\
The advantages of deliberate errors in promoting college students' memory retention and transfer & 2026 & Scopus & kvalificerar eller motsäger påståendet (0.98) \\
The benefit of retrieval practice over elaborative restudy in primary school vocabulary learning & 2014 & OA & kvalificerar eller motsäger påståendet (1.00) \\
The benefits of retrieval practice depend on item difficulty and intelligence. & 2018 & OA & kvalificerar eller motsäger påståendet (1.00) \\
The benefits of testing: Individual differences based on student factors & 2019 & OA & kvalificerar eller motsäger påståendet (0.99) \\
The costs and benefits of testing text materials & 2011 & OA & kvalificerar eller motsäger påståendet (0.98) \\
The Dark Side of Corrective Feedback: Controlled and Automatic Influences of Retrieval Practice & 2022 & Scopus & kvalificerar eller motsäger påståendet (0.98) \\
The dual-process perspective and the benefits of retrieval practice in younger and older adults & 2022 & OA & kvalificerar eller motsäger påståendet (0.96) \\
The Effect of Retrieval Practice in Primary School Vocabulary Learning & 2013 & OA & kvalificerar eller motsäger påståendet (0.99) \\
The Effectiveness of Spaced Learning, Interleaving, and Retrieval Practice in Radiology Education: A Systematic Review & 2023 & Scopus & kvalificerar eller motsäger påståendet (0.98) \\
The effectiveness of test-enhanced learning depends on trait test anxiety and working-memory capacity. & 2012 & OA & kvalificerar eller motsäger påståendet (0.98) \\
The influence of sleep-associated memory consolidation on beneficial effects of retrieval practice & 2017 & OA & kvalificerar eller motsäger påståendet (0.97) \\
The Power of Testing Memory: Basic Research and Implications for Educational Practice & 2006 & OA & kvalificerar eller motsäger påståendet (0.99) \\
The Self-Study Testing Effect with Bidirectional Paired Associate Learning & 2017 & OA & kvalificerar eller motsäger påståendet (0.98) \\
The Testing Effect and Far Transfer: The Role of Exposure to Key Information & 2016 & OA & kvalificerar eller motsäger påståendet (0.98) \\
The testing effect for mediator final test cues and related final test cues in online and laboratory experiments & 2016 & OA & kvalificerar eller motsäger påståendet (0.97) \\
The testing effect in a social setting: Does retrieval practice benefit a listener? & 2017 & OA & kvalificerar eller motsäger påståendet (0.99) \\
The Testing Effect in University Teaching: Using Multiple-Choice Testing to Promote Retention of Highly Retrievable Information & 2023 & Scopus & kvalificerar eller motsäger påståendet (0.99) \\
The testing effect with authentic educational materials: A cautionary note & 2014 & OA & kvalificerar eller motsäger påståendet (0.99) \\
The testing effect: The role of feedback and collaboration in a tertiary classroom setting & 2009 & OA & kvalificerar eller motsäger påståendet (0.99) \\
Timing's not everything: Immediate and delayed feedback are equally beneficial for performance in formative multiple-choice testing & 2024 & Scopus & kvalificerar eller motsäger påståendet (0.91) \\
Toward an Integrative Model for Managing Cognitive Load and Optimizing Retrieval Practice: A Systematic Review, Selective Meta-Analysis, and Conceptual Framework & 2025 & OA & kvalificerar eller motsäger påståendet (0.99) \\
Transitive inference and the testing effect: Retrieval practice impairs transitive inference & 2023 & Scopus & kvalificerar eller motsäger påståendet (1.00) \\
Trends in testing effect research: from lab to classroom, but not yet for all learners & 2026 & Scopus & kvalificerar eller motsäger påståendet (0.96) \\
True-False Tests Enhance Retention Relative to Rereading & 2022 & WoS & kvalificerar eller motsäger påståendet (1.00) \\
Turn-By-Turn Route Guidance Does Not Impair Route Learning & 2022 & Scopus & kvalificerar eller motsäger påståendet (0.96) \\
Untangling the benefits of multiple study opportunities and repeated testing for cued recall & 2000 & OA & kvalificerar eller motsäger påståendet (0.99) \\
Using Quizzing to Assist Student Learning in the Classroom: The Good, the Bad, and the Ugly & 2015 & Scopus & kvalificerar eller motsäger påståendet (0.99) \\
Wellness and Work: Anki, Study Habits, and Exam Outcomes in First-Year Medical Students & 2026 & Scopus & kvalificerar eller motsäger påståendet (0.76) \\
When do tests help learning? The role of working memory capacity and test frequency in secondary school students & 2026 & Scopus & kvalificerar eller motsäger påståendet (0.99) \\
When Does Practice Testing Promote Transfer on Deductive Reasoning Tasks? & 2018 & Scopus & kvalificerar eller motsäger påståendet (1.00) \\
When does retrieval induce forgetting and when does it induce facilitation? Implications for retrieval inhibition, testing effect, and text processing & 2009 & OA & kvalificerar eller motsäger påståendet (0.94) \\
Why do learners ignore expected feedback in making metacognitive decisions about retrieval practice? & 2021 & Scopus & kvalificerar eller motsäger påståendet (0.93) \\
\end{longtable}
\endgroup



\chapter{Minns man bättre det man repeterar efter ett uppehåll?}
\label{app:recap-spacing}
\chapterprecis{Författaren har ännu inte granskat resultaten i den här
  bilagan i sin helhet.}

Ja, med förbehåll: det man går igenom igen efter ett uppehåll minns man
bättre än det man övar på i ett svep, med en måttlig effekt i
klassrumsstudier, men effekten varierar mellan kurser och är för
programmering bara belagd med samband.

Den här föreläsningen repeterar hela kursen efter ett uppehåll i
programmeringen, och samma genomgång är tänkt att användas igen inför
provet.  Skälet är spridningseffekten (\foreignlanguage{english}{spacing
effect}): samma mängd övning ger bättre minne på sikt om den sprids ut i
tiden, spridd övning (\foreignlanguage{english}{distributed practice}),
än om den görs i ett svep\autocite{Cepeda2006Distributed}.  Frågan den
här bilagan besvarar är: \emph{minns man bättre det man går igenom igen
efter ett uppehåll än det man övar på i ett svep, hur stor är effekten i
verkliga kurser, och hur långt får uppehållet vara?}

\section{Metod}
\label{sec:recap-spacing-metod}

Arbetsgången är densamma som i \cref{app:recap-retrieval}, med samma
datum, verktyg och databaser och samma två databaser som saknas.  Tre
stödfrågor sökte effekten i klassrummet, i översiktsstudier och i
programmering och matematik; två motfrågor sökte uteblivna effekter och
gränser, och arbeten om hur långt uppehållet bör vara.  Alla frågor
använder samma fem namn på begreppet.  \Cref{tab:recap-spacing-queries}
visar frågorna.

% Author's note (verification and reproduction):
%   scholar session : recap-spacing (S1-S3, M1-M2 on openalex, ieee, wos,
%                     scopus, -l 40; DBLP D1-D4 answered by the anti-bot
%                     page, 0 rows)
%   searched        : 2026-10-10 (s2: key rejected)
%   classification  : as bilaga B
%   provenance      : ltnotes-repetition.bib, keys Cepeda2006Distributed
%                     Carpenter2012Spacing Latimier2021Spaced
%                     Seabrook2005Distributed Bego2024Spaced
%                     Zhang2020Distributed DonoghueHattie2021
%                     McDaniel2011Quiz (from bilaga B)
%                     Mawson2025Classroom Murray2025Mathematics
%                     GrevingRichter2019
\begin{table}[htbp]
  \centering
  \caption{Sökfrågor och antal träffar per databas, 10 oktober 2026,
    verktyget \texttt{scholar}.  S = stödfråga, M = motfråga.  D står
    för (\enquote{spacing effect} \textsc{or} \enquote{distributed
    practice} \textsc{or} \enquote{spaced practice} \textsc{or}
    \enquote{spaced retrieval} \textsc{or} \enquote{spaced learning}).
    Databaser och format som i \cref{tab:recap-retrieval-queries}.}
  \label{tab:recap-spacing-queries}
  \footnotesize
  \begin{tabular}{@{}lp{0.52\linewidth}rrrr@{}}
    \toprule
    & Nyckelord & OA & IEEE & WoS & Scopus \\
    \midrule
    S1 & D \textsc{and} (classroom \textsc{or} course \textsc{or} lecture
      \textsc{or} students)
      & 40 & 30 & 40 & 40 \\
    S2 & D \textsc{and} (\enquote{meta-analysis} \textsc{or}
      \enquote{systematic review})
      & 40 & 1 & 40 & 40 \\
    S3 & D \textsc{and} (programming \textsc{or} \enquote{computer
      science} \textsc{or} STEM \textsc{or} mathematics)
      & 40 & 34 & 39 & 40 \\
    M1 & D \textsc{and} (\enquote{no effect} \textsc{or} \enquote{null
      effect} \textsc{or} \enquote{boundary conditions} \textsc{or}
      limitations \textsc{or} \enquote{massed practice})
      & 40 & 8 & 40 & 40 \\
    M2 & D \textsc{and} (\enquote{retention interval} \textsc{or} lag
      \textsc{or} \enquote{inter-study interval} \textsc{or}
      \enquote{optimal gap})
      & 40 & 2 & 40 & 40 \\
    \bottomrule
  \end{tabular}
\end{table}

\section{Resultat}
\label{sec:recap-spacing-resultat}

\paragraph{Finns effekten?}
Ja.  Den kvantitativa sammanställningen av 317 experiment konstaterar att
studietid fördelad över flera dagar ger bättre minne på sikt än samma tid
inom några minuter, och att det uppehåll som ger bäst minne blir längre
ju längre man ska minnas: \enquote{the ISI producing maximal retention
increased as retention interval increased}
\parencite[354]{Cepeda2006Distributed}.
En översikt riktad till lärare sammanfattar samma forskning och dess
tillämpningar i undervisning\autocite{Carpenter2012Spacing}, och
metaanalysen av tio studietekniker placerar spridd övning bland de två
mest effektiva\autocite{DonoghueHattie2021}.  Kombinationen som den här
föreläsningen använder --- att plocka fram ur minnet med uppehåll
emellan --- har en egen metaanalys: spridda hämtningsövningar slog samma
övningar i ett svep med \enquote{g = 0.74}, medan det inte spelade någon
roll om uppehållen växte eller var lika långa\autocite{Latimier2021Spaced}.

\paragraph{Hur stor är effekten i verkliga kurser?}
Måttlig i genomsnitt, men varierande.  En metaanalys som bara tar med
klassrumsstudier med kursmaterial och tidsskalor som i en kurs fann
\enquote{a moderate effect in favor of distributed over massed practice
(d = 0.54, 95\% CI [0.31, 0.77])}, med en tendens till större effekter
när provet låg längre fram och för äldre
studenter\autocite{Mawson2025Classroom}.
Spridningen mellan studierna är dock mycket stor --- mer än 92 procent av
variationen beror inte på slumpen --- och små studier ger större
effekter än stora, något som författarna förklarar med de små studiernas
metoder snarare än med publiceringsbias \parencite{Mawson2025Classroom}.
I matematik är effekten mindre: \enquote{g = 0.24} för studier inom en
kurs\autocite{Murray2025Mathematics}.  Ett tidigt försök i skolan fann att
läsundervisning spridd över fler lektioner gav bättre resultat än samma
undervisning samlad\autocite{Seabrook2005Distributed}, och
klassrumsförsöket med frågornas placering fann störst vinst för frågor
som repetition inför provet\autocite{McDaniel2011Quiz}.  Men den mest
direkta prövningen i universitetskurser av den här kursens slag --- nio
inledande kurser i naturvetenskap, teknik och matematik --- fann en
signifikant effekt av spridda hämtningsövningar i bara två kurser, och
den sammanvägda effekten försvann när matematikkursen togs bort:
\enquote{The generalizability of the spacing effect across STEM courses
therefore remains unclear}\autocite{Bego2024Spaced}.  Ett förregistrerat
försök i skolan fann att utspridd omläsning av en text förlorade mot
läsning i ett svep på ett omedelbart prov och bara kom ikapp efter en
vecka\autocite{GrevingRichter2019}; det gällde omläsning, inte frågor.  Redan
sammanställningen från 2006 påpekar att studier med uppehåll i veckor
och månader, de som är relevanta för undervisning, är för få, och
återger en äldre genomgång där bara en tredjedel av studierna av
intellektuella färdigheter som huvudräkning visade någon fördel av
spridd övning \parencite[355]{Cepeda2006Distributed}.

\paragraph{Gäller det programmering?}
Det finns bara korrelationsstudier.  Bland 64 studenter som fick en
flervalsfråga om programmering om dagen och själva valde hur de övade,
fick de som spred ut övningen bättre resultat på slutprovet än de som
övade i klump\autocite{Zhang2020Distributed} --- men studenterna valde
själva, så de som spred ut övningen kan ha skilt sig på andra sätt.

\section{Diskussion och begränsningar}
\label{sec:recap-spacing-diskussion}

Att effekten finns är väl belagt; den kvantitativa sammanställningen och
metaanalysen av spridda hämtningsövningar är samstämmiga.  Hur stor den
är i en kurs är osäkrare: laboratorieförsöken mäter ofta ordlistor och
korta uppehåll, och den kursstudie som ligger närmast den här kursen
fann ett blandat resultat.  Uppehållets längd talar för föreläsningen:
eftersom det bästa uppehållet växer med hur länge man ska minnas, är ett
uppehåll på några veckor före en repetition och en ny repetition före
provet förenligt med forskningen, men ingen källa anger ett bästa
uppehåll för just det fallet.

Översikten riktad till lärare, försöket i läsundervisning,
klassrumsförsöket med frågornas placering, programmeringsstudien och
metaanalysen i matematik är lästa bara i sammanfattningen.  Metaanalysen av
spridda
hämtningsövningar är läst i sin sammanfattning och i författarnas
manuskript.

\section{Slutsats}
\label{sec:recap-spacing-slutsats}

Ja, med förbehåll: att gå igenom något igen efter ett uppehåll ger
bättre minne på sikt än att öva i ett svep, med en måttlig effekt i
klassrumsstudier, och det gäller särskilt när genomgången består av
frågor att plocka fram svaret på.  Förbehållen är att effekten i
verkliga högskolekurser varierar --- i en studie av nio kurser var den
tydlig i bara två, och i matematik är den liten --- och att belägget för
programmering bara är korrelationer.  Föreläsningen kan säga att en repetition
efter ett
uppehåll hjälper minnet, men inte hur mycket, och inte att just
uppehållets längd är den bästa.

\section{Fortsatt arbete}
\label{sec:recap-spacing-fortsatt}

Semantic Scholar och DBLP saknas, och fem källor är lästa bara i
sammanfattningen.  En studie i en programmeringsplattform om
självvald spridd övning (Zhang, Paquette och Hu 2024 i
\emph{Computers \& Education})\footnote{DOI 10.1016/j.compedu.2024.105029.}
dök upp i sökningen men kunde inte hämtas och är inte läst, och detsamma
gäller ett pågående arbete om hur uppehållet efter den första
undervisningen i programmeringsspråket MATLAB hänger ihop med
studenternas resultat\footnote{DOI 10.18260/1-2--55898; sökningen gav
ingen sammanfattning.} --- den enda träffen om uppehåll i en
programmeringskurs.

Vad fältet saknar är ett försök där uppehållet är det enda som skiljer:
samma repetition av samma innehåll, efter några dagar för en grupp och
efter några veckor för en annan, med samma slutprov i en
programmeringskurs.  Visar det en tydlig vinst för det längre
uppehållet, kan föreläsningen ange den; visar det ingen skillnad,
vilar repetitionens värde på testeffekten i \cref{app:recap-retrieval}
snarare än på uppehållet.

\needspace{10\baselineskip}
\section{Träffar som rör påståendet}

\Cref{tab:sok-recap-spacing} listar de träffar som rör påståendet.

% --- scholar LaTeX fragment --------------------------------------------
% Generated by: scholar sessions export -f table
% Session: recap-spacing
% \input this file from the body of a document whose
% preamble loads:
%   \usepackage{longtable}
%   \usepackage{booktabs}
%   \usepackage{array}
% ----------------------------------------------------------------------

\begingroup\footnotesize
\begin{longtable}{@{}%
  >{\raggedright\arraybackslash}p{0.44\textwidth}c%
  >{\raggedright\arraybackslash}p{0.13\textwidth}%
  >{\raggedright\arraybackslash}p{0.26\textwidth}@{}}
\caption{Träffar som rör påståendet, spridd övning (11 citerade, 104 som stöder påståendet och 94 som kvalificerar eller motsäger det, av 379 unika träffar; de 126 angränsande och 44 felträffarna listas inte). Skälet till varje inkludering står i sista kolumnen; för maskinklassade rader anges modellens konfidens. Databaser: OpenAlex (OA; inte open access), IEEE Xplore (IEEE), Web of Science (WoS), Scopus.}\label{tab:sok-recap-spacing}\\
\toprule
Titel & År & Leverantör & Skäl \\
\midrule
\endfirsthead
\multicolumn{4}{@{}l}{\emph{\tablename~\thetable{} (forts.)}}\\
\toprule
Titel & År & Leverantör & Skäl \\
\midrule
\endhead
\midrule \multicolumn{4}{r@{}}{\emph{forts.\ på nästa sida}}\\
\endfoot
\bottomrule
\endlastfoot
A Meta-Analysis of Ten Learning Techniques & 2021 & Scopus & \emph{citerad}: spridningseffekten \\
A Meta-Analytic Review of the Benefit of Spacing out Retrieval Practice Episodes on Retention & 2021 & Scopus & \emph{citerad}: spridningseffekten \\
A Meta-analytic Review of the Effectiveness of Spacing and Retrieval Practice for Mathematics Learning. & 2025 & OA, Scopus & \emph{citerad}: i verkliga kurser \\
Distributed and massed practice: from laboratory to classroom & 2004 & OA & \emph{citerad}: i verkliga kurser \\
Distributed Learning in the Classroom: Effects of Rereading Schedules Depend on Time of Test & 2019 & WoS & \emph{citerad}: i verkliga kurser \\
Distributed practice in verbal recall tasks: A review and quantitative synthesis. & 2006 & OA & \emph{citerad}: spridningseffekten \\
Does a Distributed Practice Strategy for Multiple Choice Questions Help Novices Learn Programming? & 2020 & WoS & \emph{citerad}: programmering \\
Single-paper meta-analyses of the effects of spaced retrieval practice in nine introductory STEM courses: is the glass half full or half empty? & 2024 & Scopus & \emph{citerad}: i verkliga kurser \\
The Distributed Practice Effect on Classroom Learning: A Meta-Analytic Review of Applied Research & 2025 & Scopus & \emph{citerad}: i verkliga kurser \\
The Effect of Testing Versus Restudy on Retention: A Meta-Analytic Review of the Testing Effect & 2014 & WoS & \emph{citerad}: testeffekten (bilaga B) \\
Using Spacing to Enhance Diverse Forms of Learning: Review of Recent Research and Implications for Instruction & 2012 & OA & \emph{citerad}: spridningseffekten \\
A dissociation between engagement and learning: Enthusiastic instructions fail to reliably improve performance on a memory task & 2017 & OA & stöder påståendet (0.93) \\
A practical application of the spacing effect to classroom learning & 1982 & OA & stöder påståendet (0.70) \\
Academic procrastination, incentivized and self-selected spaced practice, and quiz performance in an online programming problem system: An intensive longitudinal investigation & 2024 & Scopus & stöder påståendet (0.82) \\
Applying distributed practice in the schools to enhance retention of spelling words & 2023 & OA & stöder påståendet (1.00) \\
Applying evidence-based retrieval and distributed practice techniques to optimize student learning and exam scores in a university course & 2026 & WoS & stöder påståendet (0.96) \\
Applying retrieval and distributed practices to enhance student learning and achievement in a university course & 2025 & OA & stöder påståendet (0.93) \\
Applying the Spacing Effect in the Classroom & 2019 & OA & stöder påståendet (0.68) \\
Artificial grammar learning is facilitated by distributed practice: Evidence from a letter reordering task & 2022 & Scopus & stöder påståendet (0.98) \\
Benefits of spaced learning are predicted by the re-encoding of past experience in ventromedial prefrontal cortex & 2025 & Scopus & stöder påståendet (0.97) \\
Beyond Time-on-Task: The Relationship Between Spaced Study and Certification in MOOCs & 2015 & OA & stöder påståendet (0.95) \\
Chapter 14: Spaced practice in classrooms & 2024 & OA & stöder påståendet (0.87) \\
Cognitive Load Theory, Spacing Effect, and Working Memory Resources Depletion & 2019 & OA & stöder påståendet (0.91) \\
Combining flipped-classroom and spaced-repetition learning in a master-level bioinformatics course & 2025 & Scopus & stöder påståendet (0.86) \\
Comparative effectiveness of instructional design features in simulation-based education: Systematic review and meta-analysis & 2012 & OA & stöder påståendet (0.84) \\
Costs and Benefits of Spacing for Second Language Vocabulary Learning: Does Relearning Override the Positive and Negative Effects of Spacing? & 2023 & Scopus & stöder påståendet (0.99) \\
Cramming Versus Spaced Learning Induced by Strategic Design of Game-based Learning Platform & 2022 & IEEE & stöder påståendet (0.79) \\
Distributed and Massed Practices in Teaching Concepts to the Children with Autism & 2025 & WoS & stöder påståendet (0.99) \\
Distributed Practice and Distributed Representations: Investigating the Spacing Effect Using EEG & 2015 & OA & stöder påståendet (0.94) \\
Distributed Practice and Interleaved Practice: Undergraduate Students' Strategies, Experiences, and Beliefs & 2025 & OA & stöder påståendet (0.87) \\
Distributed Practice or Spacing Effect & 2020 & OA & stöder påståendet (0.99) \\
Distributing Learning Over Time: The Spacing Effect in Children's Acquisition and Generalization of Science Concepts & 2012 & OA & stöder påståendet (0.99) \\
Effect of prefrontal and parietal tDCS on learning and recognition of verbal and non-verbal material & 2016 & OA & stöder påståendet (0.91) \\
Effect of Spaced and Massed Learning Approaches on the Performance of Faculty Members in Medical Education & 2024 & WoS & stöder påståendet (0.91) \\
Effects of distributed practice on the acquisition of verb-noun collocations & 2022 & OA & stöder påståendet (0.98) \\
Effects of Massed and Distributed Practice on the Learning and Retention of Second-Language Vocabulary & 1981 & OA & stöder påståendet (1.00) \\
Effects of spaced learning on knowledge acquisition, knowledge retention, and learning-related anxiety in health professions education: A systematic review and meta-analysis of randomized controlled trials & 2026 & Scopus & stöder påståendet (0.99) \\
Effects of spacing and embellishment on memory for the main points of a text & 1982 & OA & stöder påståendet (0.87) \\
Encoding variability theory and the spacing effect in associate learning & 1976 & OA & stöder påståendet (0.90) \\
Evaluating Group Instruction and Spaced Practice to Teach Life Skills to Children with Developmental Disabilities & 2026 & Scopus & stöder påståendet (0.91) \\
Exploring the effectiveness and awareness of Anki flashcard application on medical student's academic performance: a cross-sectional study from Saudi Arabia & 2026 & WoS & stöder påståendet (0.76) \\
Extending Cognitive Load Theory to Incorporate Working Memory Resource Depletion: Evidence from the Spacing Effect & 2017 & OA & stöder påståendet (0.99) \\
Facilitating Memory for Novel Characters by Reducing Neural Repetition Suppression in the Left Fusiform Cortex & 2010 & OA & stöder påståendet (0.98) \\
Game-based Learning Approach to Cybersecurity & 2020 & IEEE & stöder påståendet (0.76) \\
Gespreide Versus Aaneengesloten Oefening in het Economisch Domein & 2018 & OA & stöder påståendet (0.84) \\
Impact of Spaced Learning on Educational Outcomes in Science Teaching & 2026 & Scopus & stöder påståendet (0.96) \\
Implementing Retrieval and Distributed Practice in a Rotating Block Schedule: Boosting Students' Memory and Metacognition & 2026 & Scopus & stöder påståendet (0.95) \\
Improvement of the modified low-first method for iterative learning & 2010 & IEEE & stöder påståendet (0.86) \\
Influence of Distributed Practice and Daily Testing on Weekly Spelling Tests 1 & 1974 & OA & stöder påståendet (0.97) \\
Investigating the effect of spaced versus massed practice on vocabulary retention in the EFL classroom & 2019 & OA & stöder påståendet (0.96) \\
Judging credibility: Can spaced lessons help students think more critically online? & 2019 & OA & stöder påståendet (0.99) \\
Lag Effects in L2 Phonetic Training: Longer Intersession Intervals Compensate for Auditory Individual Differences & 2026 & Scopus & stöder påståendet (0.98) \\
Learning Verbs in Sentences: Children With Developmental Language Disorder and the Role of Retrieval Practice & 2024 & Scopus & stöder påståendet (0.97) \\
Long-term spacing effect benefits in developmental amnesia: Case experiments in rehabilitation. & 2014 & OA & stöder påståendet (0.94) \\
Massed practice, distributed practice, and motor ability: Which one affects fencing attack skills using moving targets? & 2024 & WoS & stöder påståendet (0.82) \\
Mitigating Knowledge Decay in Engineering Calculus Through Spaced Repetition and Active Learning & 2026 & Scopus & stöder påståendet (0.99) \\
Modernizing the Flipped Classroom: Replacing Lecture Time With Asynchronous Spaced Retrieval & 2026 & Scopus & stöder påståendet (0.90) \\
More Steps, Same Effect: Spacing Increases the Retention of Mathematics Procedures of Varying Complexity & 2024 & OA & stöder påståendet (0.99) \\
Neural mechanisms of the spacing effect in episodic memory: A parallel EEG and fMRI study & 2015 & OA & stöder påståendet (0.98) \\
On effective learning of English collocations: From perspectives of distributed practice and semantic restructuring & 2023 & OA & stöder påståendet (0.88) \\
On the Acceptance and Effects of Recapping Self-Test Questions in MOOCs & 2020 & IEEE & stöder påståendet (0.77) \\
Optimising learning using flashcards: Spacing is more effective than cramming & 2009 & OA & stöder påståendet (1.00) \\
Optimizing distributed practice online & 2025 & OA & stöder påståendet (1.00) \\
Optimizing L2 phonetic learning & 2026 & Scopus & stöder påståendet (0.99) \\
Optimizing song retention through the spacing effect & 2021 & Scopus & stöder påståendet (0.99) \\
People mind wander more during massed than spaced inductive learning. & 2015 & OA & stöder påståendet (0.97) \\
Practice and Forgetting Effects on Vocabulary Memory: An ActivationBased Model of the Spacing Effect & 2005 & OA & stöder påståendet (0.99) \\
PREDICTING AND IMPROVING MEMORY RETENTION Psychological Theory Matters in the Big Data Era & 2017 & WoS & stöder påståendet (0.91) \\
Preservation of long-term memory in older adults using a spaced learning paradigm & 2023 & Scopus & stöder påståendet (0.98) \\
Realization of an Effective Spaced Learning Schedule Based on a Reactivation Theory of the Spacing Effect & 1998 & OA & stöder påståendet (0.95) \\
Retrieval practice and spacing effects in multi-session treatment of naming impairment in aphasia & 2019 & Scopus & stöder påståendet (0.94) \\
Retrieval practice and spacing in an engineering mathematics classroom: Do the effects add up? & 2017 & IEEE & stöder påståendet (1.00) \\
Retrieval Practice Through an Integrated Anatomy Laboratory Experience in a Doctor of Physical Therapy Curriculum & 2026 & WoS & stöder påståendet (0.79) \\
Revolutionizing radiology education: exploring and implementing spaced repetition in radiology teaching and curricula & 2026 & Scopus & stöder påståendet (0.90) \\
Role of the hippocampus in the spacing effect during memory retrieval & 2020 & Scopus & stöder påståendet (0.98) \\
Salient Memory: Effects of Distributed Learning on Cortical Regions During Memory Retrieval & 2025 & Scopus & stöder påståendet (0.98) \\
Semantic-feature variability and the spacing effect & 1979 & OA & stöder påståendet (0.83) \\
Spaced example-based learning facilitates advanced mathematics learning in college students & 2024 & Scopus & stöder påståendet (0.99) \\
Spaced instruction for psychological concepts: Insights from interpersonal brain synchronization and students' perceptions of teaching & 2026 & Scopus & stöder påståendet (0.99) \\
Spaced Learning Enhances Episodic Memory by Increasing Neural Pattern Similarity Across Repetitions & 2019 & WoS & stöder påståendet (0.98) \\
Spaced Learning: Its Implications in the Language Classroom & 2014 & OA & stöder påståendet (0.91) \\
Spaced repetition and active recall improves academic performance among pharmacy students & 2026 & Scopus & stöder påståendet (0.82) \\
Spaced Repetition and Retrieval Practice: Efficient Learning Mechanisms from a Cognitive Psychology Perspective and Their Empowerment by AI & 2025 & OA & stöder påståendet (0.86) \\
Spaced Repetition Learning in Radiology Education: Exploring Its Potential and Practical Application & 2025 & Scopus & stöder påståendet (0.88) \\
Spaced Retrieval Practice in Two Fundamental Engineering Courses: Calculus and Physics & 2022 & WoS & stöder påståendet (0.87) \\
Spacing as the friend of both memory and induction in young and older adults. & 2010 & OA & stöder påståendet (0.99) \\
Spacing effects in learning and memory & 2025 & Scopus & stöder påståendet (0.92) \\
Spacing effects in realworld classroom vocabulary learning & 2010 & OA & stöder påståendet (0.99) \\
Spacing Simultaneously Promotes Multiple Forms of Learning in Children's Science Curriculum & 2014 & OA & stöder påståendet (1.00) \\
Specifying the neural basis of the spacing effect with multivariate ERP & 2020 & Scopus & stöder påståendet (0.97) \\
Spreading the words: A spacing effect in vocabulary learning & 2012 & OA & stöder påståendet (1.00) \\
STEMing evidence-based study strategies: a systematic review of literature & 2026 & Scopus & stöder påståendet (0.97) \\
SYSTEMATIC REVIEW ON LEARNING CONSOLIDATION IN UNIVERSITY STUDENTS & 2025 & WoS & stöder påståendet (0.79) \\
Teach Students to Study Using Quizzes, Study Behavior Visualization, and Reflection: A Case Study in an Introduction to Programming Course & 2023 & Scopus & stöder påståendet (0.79) \\
Testing the reminding account of the lag effect in L2 vocabulary learning & 2022 & WoS & stöder påståendet (0.98) \\
The Breakfast Club: Enhancing Emergency Medicine Education Through Spaced Retrieval and Elaborative Interrogation Techniques & 2026 & WoS & stöder påståendet (0.82) \\
The distributed practice effects of speaking task repetition & 2021 & OA & stöder påståendet (0.98) \\
The Effect of Context on Distributed Practice & 2013 & OA & stöder påståendet (0.92) \\
The Effect of Weekly Distributed Mathematics Homework and Quizzes on the Learning Performance of Engineering Students & 2019 & WoS & stöder påståendet (0.78) \\
The Effectiveness of Spaced Repetition in Medical Education: A Systematic Review and Meta-Analysis & 2026 & Scopus & stöder påståendet (0.99) \\
The Effects of Repetition on Incidental Vocabulary Learning: A Meta-Analysis of Correlational Studies & 2019 & Scopus & stöder påståendet (0.86) \\
The Impact of a Blended Learning Approach Incorporating Microlearning and Spaced Learning on Clinical Competence in Nursing Students & 2025 & WoS & stöder påståendet (0.95) \\
The impact of distributed practice schedules within a classwide computation intervention: A randomized control trial & 2026 & Scopus & stöder påståendet (0.94) \\
The Impact of Spaced Learning Strategies on Vocabulary Retention in Primary English Classrooms & 2024 & OA & stöder påståendet (0.93) \\
The retention interval model & 1981 & OA & stöder påståendet (0.94) \\
The situation with respect to the spacing of repetitions and memory & 1970 & OA & stöder påståendet (0.91) \\
The Spacing Effect in Children's Generalization of Knowledge: Allowing Children Time to Forget Promotes Their Ability to Learn & 2014 & OA & stöder påståendet (0.98) \\
The spacing effect in the learning of word pairs & 1974 & OA & stöder påståendet (0.90) \\
The spacing effect stands up to big data & 2019 & WoS & stöder påståendet (0.99) \\
Unlocking words and fluency: Spaced Retrieval small-scale case study practice with Spanish-speaking A1 EFL adult learners & 2026 & Scopus & stöder påståendet (0.86) \\
Unveiling the Art of Effective Learning through Spaced Repetition and Evidence-Based Techniques & 2024 & IEEE & stöder påståendet (0.86) \\
Using an Educational Mobile Application for Learning the Essence 1.2 Kernel Alphas & 2021 & Scopus & stöder påståendet (0.70) \\
Using Spaced Learning for Secondary Science Revision: The SMART Spaces Programme & 2019 & OA & stöder påståendet (0.91) \\
When do tests help learning? The role of working memory capacity and test frequency in secondary school students & 2026 & WoS & stöder påståendet (0.96) \\
Working memory capacity and the spacing effect in cued recall & 2017 & OA & stöder påståendet (0.98) \\
A Meta-analysis of the Spacing Effect in Verbal Learning: Implications for Research on Advertising Repetition and Consumer Memory & 2003 & OA & kvalificerar eller motsäger påståendet (0.96) \\
A pilot randomized controlled trial comparing the effectiveness of different spaced learning models used during school examination revision: the SMART Spaces 24/10 model & 2023 & Scopus & kvalificerar eller motsäger påståendet (0.99) \\
A Review of the Application of Distributed Practice Principles to Naming Treatment in Aphasia & 2020 & Scopus & kvalificerar eller motsäger påståendet (0.93) \\
About Practice & 2013 & OA & kvalificerar eller motsäger påståendet (0.98) \\
Age-related differences in the impact of spacing, lag, and retention interval. & 1989 & OA & kvalificerar eller motsäger påståendet (0.98) \\
Auditory Training for Adults Who Have Hearing Loss: A Comparison of Spaced Versus Massed Practice Schedules & 2017 & OA & kvalificerar eller motsäger påståendet (0.98) \\
Avoiding Surgical Skill Decay: A Systematic Review on the Spacing of Training Sessions & 2017 & OA & kvalificerar eller motsäger påståendet (0.94) \\
Between-list lag effects in recall depend on retention interval & 2014 & OA & kvalificerar eller motsäger påståendet (1.00) \\
Beyond the Distributed Practice Effect: Is Distributed Learning Also Effective for Learning With Non-repeated Text Materials? & 2021 & Scopus & kvalificerar eller motsäger påståendet (0.99) \\
BOARD \# 81: WIP: Student outcomes as related to the interval between initial MATLAB instruction, potential interim programming encounters, and an intermediate MATLAB course & 2025 & Scopus & kvalificerar eller motsäger påståendet (0.57) \\
Cognitive fatigue induced by spaced and massed practice: effects on idiom accuracy across L2 proficiency levels & 2026 & Scopus & kvalificerar eller motsäger påståendet (0.96) \\
Conceptualising spaced learning in health professions education: A scoping review & 2020 & Scopus & kvalificerar eller motsäger påståendet (0.88) \\
Diminished But Not Forgotten: Effects of Aging on Magnitude of Spacing Effect Benefits & 2012 & OA & kvalificerar eller motsäger påståendet (1.00) \\
Distributed practice and time pressure interact to affect learning and retention of arithmetic facts & 2023 & OA & kvalificerar eller motsäger påståendet (0.99) \\
Distributing mathematical practice of third and seventh graders: A pplicability of the spacing effect in the classroom & 2018 & OA & kvalificerar eller motsäger påståendet (1.00) \\
Distribution of grammar learning through writing tasks: A potential role of proficiency & 2024 & Scopus & kvalificerar eller motsäger påståendet (0.98) \\
Distribution of Practice Effects in the Acquisition and Retention of L2 Mandarin Tonal Word Production & 2019 & WoS & kvalificerar eller motsäger påståendet (0.98) \\
Do worked examples boost the spacing effect on lasting learning? & 2025 & Scopus & kvalificerar eller motsäger påståendet (0.97) \\
Does automated reassessment and relearning improve mathematics performance better than computer adaptive practice alone? & 2025 & Scopus & kvalificerar eller motsäger påståendet (0.99) \\
Does Learning Method Matter in Cyber Security Behaviour? Spaced Vs. Massed e-Learning & 2021 & IEEE & kvalificerar eller motsäger påståendet (0.96) \\
Durable Learning Strategies in Nursing Education: State-of-the-Evidence Review & 2024 & Scopus & kvalificerar eller motsäger påståendet (0.91) \\
Effects of distributed practice and criterion level on word retrieval in aphasia & 2020 & Scopus & kvalificerar eller motsäger påståendet (0.91) \\
Effects of massed and spaced task repetitions on L2 writing task performance and task emotions & 2026 & Scopus & kvalificerar eller motsäger påståendet (0.98) \\
Effects of spacing on contextual vocabulary learning: Spacing facilitates the acquisition of explicit, but not tacit, vocabulary knowledge & 2020 & OA & kvalificerar eller motsäger påståendet (1.00) \\
Effects of spacing on the acquisition of explicit and implicit vocabulary knowledge: An approximate replication of Nakata and Elgort (2021) & 2026 & Scopus & kvalificerar eller motsäger påståendet (0.99) \\
Effects of the 4/3/2 activity revisited: Aptitude matters, but distributed practice does not & 2026 & Scopus & kvalificerar eller motsäger påståendet (0.99) \\
Evaluating the Spacing Effect Theory on the Instructional Effectiveness of Semester-Length versus Quarter-Length Introductory Computer Literacy Courses in Institutions of Higher Learning & 2013 & OA & kvalificerar eller motsäger påståendet (0.75) \\
Examining the benefits of spaced retrieval practice on transfer and the student experience in the classroom & 2026 & Scopus & kvalificerar eller motsäger påståendet (0.97) \\
Examining the effects of training on young and older adult implementation of spaced retrieval strategies & 2022 & Scopus & kvalificerar eller motsäger påståendet (0.96) \\
Expanded vs. Equal Interval Spaced Retrieval Practice: Exploring Different Schedules of Spacing and Retention Interval in Younger and Older Adults & 2008 & OA & kvalificerar eller motsäger påståendet (0.98) \\
Exploring Distribution-of-Practice Effects on L2 Grammar Learning Among Older Adults & 2026 & Scopus & kvalificerar eller motsäger påståendet (0.98) \\
Free Recall of Items Presented after Massed- and Distributed-Practice Items & 1972 & OA & kvalificerar eller motsäger påståendet (0.97) \\
Frequency judgments and the spacing effect: Immediate and delayed performance & 1976 & OA & kvalificerar eller motsäger påståendet (0.76) \\
Identifying the optimal thoracentesis training strategy: a randomized non-inferiority study & 2025 & WoS & kvalificerar eller motsäger påståendet (0.96) \\
Initial Practice Performance Moderates the Distributed Practice Effect in Complex Procedural Knowledge & 2025 & WoS & kvalificerar eller motsäger påståendet (0.99) \\
Input spacing and the learning of L2 vocabulary in a classroom context & 2020 & WoS & kvalificerar eller motsäger påståendet (0.99) \\
Judging the Credibility of Websites: An Effectiveness Trial of the Spacing Effect in the Elementary Classroom & 2019 & OA & kvalificerar eller motsäger påståendet (0.88) \\
Judging the credibility of websites: an effectiveness trial of the spacing effect in the elementary classroom & 2022 & OA & kvalificerar eller motsäger påståendet (0.99) \\
Lack of spacing effects during piano learning & 2017 & OA & kvalificerar eller motsäger påståendet (0.95) \\
Lag effects for foreign language vocabulary learning through Quizlet & 2024 & WoS & kvalificerar eller motsäger påståendet (0.99) \\
Lag effects in grammar learning: A desirable difficulties perspective & 2022 & WoS & kvalificerar eller motsäger påståendet (0.99) \\
Learning techniques that medical students use for long-term retention: A cross-sectional analysis & 2023 & Scopus & kvalificerar eller motsäger påståendet (0.83) \\
Learning techniques to improve memory in children: a systematic review & 2025 & Scopus & kvalificerar eller motsäger påståendet (0.99) \\
Level of initial training moderates the effects of distributing practice over multiple days with expanding, contracting, and uniform schedules: Evidence for study-phase retrieval & 2018 & OA & kvalificerar eller motsäger påståendet (0.98) \\
No Robust Effect of Distributed Practice on the Short- and Long-Term Retention of Mathematical Procedures & 2020 & Scopus & kvalificerar eller motsäger påståendet (1.00) \\
On the Effects of Spaced Versus Massed Practice on Deliberate Learning of English Idioms across L2 Proficiency Levels & 2026 & Scopus & kvalificerar eller motsäger påståendet (1.00) \\
Optimizing Distributed Practice & 2009 & OA & kvalificerar eller motsäger påståendet (0.99) \\
Parallels between spacing effects during behavioral and cellular learning & 2012 & OA & kvalificerar eller motsäger påståendet (0.94) \\
Rapid Cycle Deliberate Practice With Spaced Training Achieves Comparable Results in Achieving Simple Laparoscopic Skills to Traditional Weekly Training and Saves Time & 2025 & WoS & kvalificerar eller motsäger påståendet (0.95) \\
Spaced Digital Education for Health Professionals: Systematic Review and Meta-Analysis & 2024 & Scopus & kvalificerar eller motsäger påståendet (0.99) \\
Spaced learning and melodic memory & 2025 & Scopus & kvalificerar eller motsäger påståendet (0.99) \\
Spaced Learning Supports Productive Struggle in an Online Learning Platform & 2026 & Scopus & kvalificerar eller motsäger påståendet (0.99) \\
Spaced learning versus massed learning in resuscitation --- A systematic review & 2020 & Scopus & kvalificerar eller motsäger påståendet (0.99) \\
Spaced Retrieval Practice Improves Engineering Student Performance in Physics & 2023 & IEEE & kvalificerar eller motsäger påståendet (1.00) \\
Spacing and Interleaving Effects Require Distinct Theoretical Bases: a Systematic Review Testing the Cognitive Load and Discriminative-Contrast Hypotheses & 2021 & Scopus & kvalificerar eller motsäger påståendet (0.99) \\
Spacing and task complexity in mathematics learning & 2025 & OA & kvalificerar eller motsäger påståendet (0.99) \\
Spacing and Testing Effects & 2010 & OA & kvalificerar eller motsäger påståendet (0.99) \\
Spacing Effects in Learning & 2008 & OA & kvalificerar eller motsäger påståendet (1.00) \\
Spacing Effects in Task Repetition Research & 2023 & Scopus & kvalificerar eller motsäger påståendet (0.94) \\
Spacing learning units affects both learning and forgetting & 2022 & Scopus & kvalificerar eller motsäger påståendet (0.99) \\
Spacing repetitions over 1 week & 1980 & OA & kvalificerar eller motsäger påståendet (0.99) \\
Systematic review of distributed practice and retrieval practice in health professions education & 2024 & Scopus & kvalificerar eller motsäger påståendet (0.99) \\
Systematic Review of the Effectiveness of Neurodidactic Spaced Learning Strategies in Long-Term Memory & 2026 & Scopus & kvalificerar eller motsäger påståendet (0.99) \\
Ten years of massed practice on distributed practice. & 1961 & OA & kvalificerar eller motsäger påståendet (0.93) \\
The Cognitive Mechanisms of the Positivity Reactivity Effect on Word Recognition Memory & 2026 & Scopus & kvalificerar eller motsäger påståendet (0.88) \\
The design and optimization of SMART Spaces: A science revision program utilizing a spaced learning framework. & 2018 & OA & kvalificerar eller motsäger påståendet (0.87) \\
The difficulty to evoke the spacing effect in mathematics: New findings and theoretical considerations & 2025 & OA & kvalificerar eller motsäger påståendet (0.91) \\
The Effect of Practice Distribution on Skill Retention in Virtual Reality Temporal Bone Surgery Training & 2019 & IEEE & kvalificerar eller motsäger påståendet (0.96) \\
The effectiveness of brief, spaced practice on student difficulties with basic and essential engineering skills & 2013 & IEEE & kvalificerar eller motsäger påståendet (0.80) \\
The Effectiveness of Spaced Learning, Interleaving, and Retrieval Practice in Radiology Education: A Systematic Review & 2023 & Scopus & kvalificerar eller motsäger påståendet (0.98) \\
The effects of distributed practice on second language fluency development & 2024 & OA & kvalificerar eller motsäger påståendet (0.97) \\
The Effects of Practice Schedules on the Acquisition and Retention of a Grapho-Motor Skill in Young-Adults & 2025 & Scopus & kvalificerar eller motsäger påståendet (0.90) \\
The Effects of Spaced and Massed Practice Schedule on L2 Contextualized Grammar Learning of a Simple and a Complex Structure & 2026 & Scopus & kvalificerar eller motsäger påståendet (0.95) \\
The Effects of Spaced Practice on Second Language Learning: A Meta-Analysis & 2022 & Scopus & kvalificerar eller motsäger påståendet (1.00) \\
The Efficacy of Massed Versus Distributed Practice As a Function of Desired Learning Outcomes and Grade Level of the Student & 2021 & OA & kvalificerar eller motsäger påståendet (0.99) \\
The Impact of Simulation-Based Spaced Training for Skills Acquisition on Learning and Performance Outcomes Among Healthcare Professionals & 2026 & Scopus & kvalificerar eller motsäger påståendet (0.99) \\
The Impact of Spaced and Massed Practice on Deliberate Learning of English Opaque Idioms across L2 Proficiency Levels: A Mixed-Methods Research & 2024 & OA & kvalificerar eller motsäger påståendet (0.98) \\
The influence of the spacing effect on L2 vocabulary learning: A study on Chinese university students & 2023 & OA & kvalificerar eller motsäger påståendet (1.00) \\
The L2 proficiency puzzle: spaced and massed practice in deliberate learning of English collocations & 2026 & WoS & kvalificerar eller motsäger påståendet (0.98) \\
The role of lag effect in distributed practice on learning novel vocabulary & 2023 & Scopus & kvalificerar eller motsäger påståendet (0.91) \\
The Spacing Effect in Aircraft Recognition & 1996 & OA & kvalificerar eller motsäger påståendet (0.96) \\
The spacing effect: Implications for relearning & 1997 & OA & kvalificerar eller motsäger påståendet (0.99) \\
The time between tasks in task repetition research: A systematic review & 2025 & Scopus & kvalificerar eller motsäger påståendet (0.92) \\
Theoretical Implications of the Spacing Effect 1 & 2024 & OA & kvalificerar eller motsäger påståendet (0.97) \\
Time distribution and intentional vocabulary learning through repeated reading: a partial replication and extension & 2023 & Scopus & kvalificerar eller motsäger påståendet (0.98) \\
Translating Laboratory Evidence into Classroom Practice with Teacher-Led Randomized Controlled Trials---A Perspective and Meta-Analysis & 2020 & Scopus & kvalificerar eller motsäger påståendet (0.91) \\
U-Behavior: a method to promote spaced and interleaved retrieval practice in challenging undergraduate microbiology and biomedical sciences courses & 2026 & Scopus & kvalificerar eller motsäger påståendet (0.94) \\
Understanding the underlying mechanism of the spacing effect in verbal learning: a case for encoding variability and study-phase retrieval & 2016 & OA & kvalificerar eller motsäger påståendet (0.98) \\
Variability in knowledge retention of medical students: repeated and recently learned basic science topics & 2025 & Scopus & kvalificerar eller motsäger påståendet (0.67) \\
Verbal and spatial acquisition as a function of distributed practice and code-specific interference & 2019 & OA & kvalificerar eller motsäger påståendet (0.95) \\
Very Similar Spacing-Effect Patterns in Very Different Learning/Practice Domains & 2014 & OA & kvalificerar eller motsäger påståendet (0.99) \\
What is the role of affect in the spacing effect? & 1983 & OA & kvalificerar eller motsäger påståendet (0.97) \\
Why Desirable Difficulties `Work': A Review of the Evidence From Cognitive and Educational Psychology and Some Caveats for the Health Professions Education Field & 2026 & Scopus & kvalificerar eller motsäger påståendet (0.95) \\
Word learning by children with developmental language disorder: Identifying gaps in our understanding of spaced retrieval effects & 2024 & OA & kvalificerar eller motsäger påståendet (0.97) \\
\end{longtable}
\endgroup



\chapter{Lär man sig mer om man får frågor innan något undervisas?}
\label{app:recap-pretest}
\chapterprecis{Författaren har ännu inte granskat resultaten i den här
  bilagan i sin helhet.}

Ja, men smalt: en fråga om något innan det undervisas hjälper en att lära
sig just det som frågan gällde, när svaret kommer efteråt, men nästan
inte alls annat; i verkliga kurser är vinsten liten och ojämn, och den
avtar efter en till två veckor.

Den här föreläsningen ger också en försmak av det som ännu inte
undervisats, med frågor om det som ska komma: avsnitten från
\cref{sec:recap-mangder} och framåt.
Tanken är att en fråga som ställs innan svaret har lärts ut --- ett
förtest (\foreignlanguage{english}{pretesting}) eller en förfråga
(\foreignlanguage{english}{prequestion}) --- gör att man lär sig mer när
svaret sedan kommer, även om man svarade fel\autocite{Richland2009Pretesting},
och att en översikt i förväg, en förhandsstruktur
(\foreignlanguage{english}{advance organizer}), gör det nya lättare att
fästa vid det man redan kan\autocite{Luiten1980Organizers}.  Frågan den
här bilagan besvarar är: \emph{hjälper det att få frågor om, eller en
översikt över, något innan det undervisas --- och i så fall för vad?}

\section{Metod}
\label{sec:recap-pretest-metod}

Arbetsgången är densamma som i \cref{app:recap-retrieval}, med samma
datum, verktyg och databaser.  Två stödfrågor sökte förtest och
förfrågor, i allmänhet och i klassrum och föreläsningar; en tredje sökte
förhandsstrukturer och översikter.  Två motfrågor sökte uteblivna
effekter och gränser för vart och ett av de två begreppen, med samma
namn på begreppen som stödfrågorna.  \Cref{tab:recap-pretest-queries}
visar frågorna.  IEEE Xplore gav inga träffar på förtest, vilket
stämmer med att begreppet hör till psykologin och inte till teknik.

% Author's note (verification and reproduction):
%   scholar session : recap-pretest (S1-S3, M1-M2 on openalex, ieee, wos,
%                     scopus, -l 40; DBLP D1-D3 answered by the anti-bot
%                     page, 0 rows)
%   searched        : 2026-10-10 (s2: key rejected)
%   classification  : as bilaga B
%   provenance      : ltnotes-repetition.bib, keys Richland2009Pretesting
%                     PanCarpenter2023Prequestioning StHilaire2024Guessing
%                     Carpenter2018Prequestions Geller2017Prequestions
%                     Toftness2018Prequestions Pan2020MindWandering
%                     PanSana2021Pretesting Latimier2019Pretesting
%                     JamesStorm2019Pretesting Kliegl2024Pretesting
%                     Luiten1980Organizers Stone1983Organizers
%                     Mayer1979Organizers KingShepard2025Prequestions
%                     Motz2026ManyClasses Pan2026Temporal
%                     SanaCarpenter2023Placement; Yang2021Testing (bilaga B);
%                     SinhaKapur2021 reused from Övning: Funktioner och
%                     variabler, bilaga B (block copied)
\begin{table}[htbp]
  \centering
  \caption{Sökfrågor och antal träffar per databas, 10 oktober 2026,
    verktyget \texttt{scholar}.  S = stödfråga, M = motfråga.  P står
    för (\enquote{pretesting effect} \textsc{or} \enquote{pretest
    effect} \textsc{or} prequestion \textsc{or} prequestions \textsc{or}
    \enquote{errorful generation} \textsc{or} \enquote{pre-questions}),
    O för (\enquote{advance organizer} \textsc{or} \enquote{advance
    organizers} \textsc{or} preview \textsc{or} previewing).  Databaser
    och format som i \cref{tab:recap-retrieval-queries}.}
  \label{tab:recap-pretest-queries}
  \footnotesize
  \begin{tabular}{@{}lp{0.52\linewidth}rrrr@{}}
    \toprule
    & Nyckelord & OA & IEEE & WoS & Scopus \\
    \midrule
    S1 & P \textsc{and} (learning \textsc{or} retention \textsc{or}
      memory)
      & 40 & 0 & 40 & 40 \\
    S2 & P \textsc{and} (classroom \textsc{or} lecture \textsc{or} video
      \textsc{or} course \textsc{or} students)
      & 40 & 0 & 40 & 40 \\
    S3 & O \textsc{and} (lecture \textsc{or} instruction \textsc{or}
      learning) \textsc{and} (experiment \textsc{or}
      \enquote{meta-analysis} \textsc{or} effect)
      & 40 & 40 & 40 & 40 \\
    M1 & P \textsc{and} (\enquote{non-prequestioned} \textsc{or}
      \enquote{no benefit} \textsc{or} \enquote{null effect} \textsc{or}
      \enquote{boundary conditions} \textsc{or} limitations \textsc{or}
      \enquote{does not})
      & 40 & 0 & 12 & 14 \\
    M2 & O \textsc{and} (lecture \textsc{or} instruction \textsc{or}
      learning) \textsc{and} (\enquote{no effect} \textsc{or} ineffective
      \textsc{or} criticism \textsc{or} \enquote{no significant})
      & 40 & 1 & 34 & 40 \\
    \bottomrule
  \end{tabular}
\end{table}

\section{Resultat}
\label{sec:recap-pretest-resultat}

\paragraph{Hjälper en fråga innan svaret har lärts ut?}
Ja, för just det som frågan gällde.  I det experiment som gav effekten
dess namn svarade deltagarna på frågor om en text innan de läste den,
och de mindes mer än de som fick läsa längre --- i alla fem experimenten,
och även när bara de frågor räknades som de hade svarat fel
på \parencite[243]{Richland2009Pretesting}.  En aktuell översikt
sammanfattar att förfrågor och förtest \enquote{can often enhance
learning}, förutsatt att man får studera de rätta svaren efteråt, men
att vinsten varierar med upplägg och med hur lärandet
mäts\autocite{PanCarpenter2023Prequestioning}.  Effekten finns också med
inspelade föreläsningar, där förtest minskade tankeflykten och gav
bättre resultat\autocite{Pan2020MindWandering}, och den kan växa med
tiden: i ett försök med en text fördubblades den ungefär när provet
flyttades till en vecka senare\autocite{Kliegl2024Pretesting}.

\paragraph{Hjälper den för det som frågan inte gällde?}
Nej, i stort sett inte, och det är bilagans viktigaste förbehåll.  En
förregistrerad metaanalys skiljer de två: förfrågor har \enquote{a
moderate specific effect ( g = 0.54, k = 97) but a virtually nonexistent
general effect ( g = 0.04, k = 91)}\autocite{StHilaire2024Guessing},
och en andra, oberoende metaanalys fann samma mönster, \enquote{g = .66}
mot \enquote{g = .01}, och en starkare effekt när förfrågorna följdes
av återkoppling\autocite{KingShepard2025Prequestions}.
Fem experiment fann ingen effekt alls på det som inte
förtestades\autocite{JamesStorm2019Pretesting},
och med verkliga föreläsningsvideor på över 20 minuter hjälpte
förfrågorna det som frågades om men inte, eller bara marginellt,
resten\autocite{Toftness2018Prequestions}.  Ett undantag är ordningen:
det som inte förtestades gynnades när det kom före det förtestade i
lektionen, inte efter\autocite{SanaCarpenter2023Placement}.

\paragraph{Hjälper det i en verklig kurs?}
Lite, enligt den största studien, och de mindre klassrumsstudierna är
blandade.  I 30 klasser från skolans årskurs sex till universitetets
sista år, med 1\,571 elever och studenter, gav förfrågor före en
inspelad föreläsning bättre resultat på senare prov som innehöll
förfrågorna, men vinsten var liten: 73,7 mot 68,8 procent rätt, och
författarna påpekar att en effekt av laboratoriets storlek
\enquote{would be far-fetched in a real classroom
context}\autocite{Motz2026ManyClasses}.  En del studenter hoppade dessutom
över föreläsningen helt, och då gav förfrågorna ingen vinst alls; vinsten
var störst för dem som svarade rätt, vilket författarna beskriver som
\enquote{a rich-get-richer scenario} \parencite{Motz2026ManyClasses}.  I en
kurs i psykologi hjälpte förfrågorna
bara för det som frågades om, och en vecka senare lade de lite eller
inget till vinsten av frågor efter
undervisningen\autocite{Carpenter2018Prequestions}.
I en kurs i kemiteknik gav frågor efter lektionen en tydlig vinst, men
förfrågorna i början av lektionen \enquote{do not appear to enhance
learning, nor to enhance the effects of retrieval
practice}\autocite{Geller2017Prequestions}.  Klassrumsmetaanalysen i
\cref{app:recap-retrieval} fann en liten men signifikant vinst för frågor
före lektionen, \enquote{g = 0.186}, mot \enquote{g = 0.536} för frågor
efter den \parencite{Yang2021Testing}.  I en kurs där studenterna fick
korta frågor med svarsapparat innan begreppen undervisades, mindes de
begreppen bättre på provet veckor senare, men vinsten nådde inte
konsekvent frågor som krävde förståelse \parencite{Pan2019Jargon}.

\paragraph{Hur länge håller effekten?}
Under en till två veckor väl, därefter sämre.  I ett försök med en
text ungefär fördubblades effekten när provet flyttades från direkt
efteråt till en vecka senare \parencite{Kliegl2024Pretesting}, och i
försök med uppehåll upp till 30 dagar var vinsten tydligast under den
första veckan, märkbar i upp till två veckor och avtagande
därefter\autocite{Pan2026Temporal}.  Uppehållet räknas där från förtestet
och undervisningen till provet, inte från förtestet till undervisningen.

\paragraph{Är förtest bättre eller sämre än frågor efteråt?}
Källorna är oeniga.  I fem experiment med faktatexter gav förtest
högre resultat än frågor efter läsningen\autocite{PanSana2021Pretesting};
i ett annat experiment med prov efter en vecka gav båda en vinst, men
frågor efteråt en dubbelt så stor (\enquote{d = 0.74} mot \enquote{d =
0.35}), och bara de förde över till nya
frågor\autocite{Latimier2019Pretesting}.

\paragraph{Hjälper en översikt i förväg?}
Ja, enligt två äldre metaanalyser, men med en omstridd historia.  En
genomgång av 135 studier fann att förhandsstrukturer underlättar både
lärande och minne \parencite{Luiten1980Organizers}, och en andra med 112
studier fann detsamma men inte i det mönster teorin bakom dem
förutsäger\autocite{Stone1983Organizers}.  En tidigare översikt hade
kommit fram till att de inte hjälper alls; Mayer menar att den
översikten hade brister och att förhandsstrukturer hjälper \enquote{if
used in appropriate situations and measured
properly}\autocite{Mayer1979Organizers}.  Ingen av sammanfattningarna
anger hur stor effekten är.

\paragraph{Ett närliggande belägg.}
Att först försöka lösa ett problem och sedan få undervisning om det ---
produktivt misslyckande --- ger måttligt bättre förståelse än omvänd
ordning, men effekten vänder för yngre elever och för
domänöverskridande färdigheter\autocite{SinhaKapur2021}; det
vetenskapliga underlaget för det finns i \emph{Övning: Funktioner och
variabler}, bilaga~B.  Det gäller problem att lösa, inte korta frågor,
men pekar åt samma håll: ett misslyckat försök före undervisningen är
inte bortkastat.

\section{Diskussion och begränsningar}
\label{sec:recap-pretest-diskussion}

Effekten på det som frågades om är väl belagd, och att den inte sprider
sig till resten är lika väl belagt.  För den här föreläsningen betyder
det att frågorna om det som ska komma hjälper för de begrepp de handlar
om, inte för området i stort.  I verkliga kurser är resultaten blandade:
den största studien fann en liten vinst i genomsnitt, omkring fem
procentenheter, medan två mindre
klassrumsstudier inte fann någon vinst utöver frågorna efter
undervisningen.  Ingen källa gäller programmering, och ingen prövar en
försmak som ges veckor innan undervisningen, vilket är fallet här: i
källorna kommer undervisningen strax efter förfrågan, och uppehållet
fram till provet räcker som längst 30 dagar, med avtagande effekt efter
en till två veckor.

Många av källorna är lästa bara i sammanfattningen: klassrumsstudien i
psykologi, studien med föreläsningsvideor, studien om tankeflykt,
studien om det som inte förtestades, de två studierna med fördröjt prov, den
andra metaanalysen av förfrågor,
de två metaanalyserna av förhandsstrukturer och Mayers översikt.  Den äldre
negativa översikten av förhandsstrukturer är inte läst; den är återgiven
genom Mayer.

\section{Slutsats}
\label{sec:recap-pretest-slutsats}

Ja, men smalt: en fråga om något innan det undervisas hjälper en att
lära sig just det som frågan gällde, när svaret kommer efteråt, men
nästan inte alls det som frågan inte gällde; i verkliga kurser är
vinsten liten (omkring fem procentenheter i den största studien) och
ojämnt fördelad, och effekten avtar
efter en till två veckor.  Föreläsningen kan därför säga att frågorna om
det som ska komma förbereder de begrepp de handlar om, men inte att en
försmak förbereder hela området.  Frågorna bör handla om de begrepp som
är viktigast i det som kommer, följas av svaret, och helst ställas igen
när begreppet undervisas.

\section{Fortsatt arbete}
\label{sec:recap-pretest-fortsatt}

Semantic Scholar och DBLP saknas, och tio källor är lästa bara i
sammanfattningen (se \cref{sec:recap-pretest-diskussion}).  Översikten
gav inget svar på hur stor effekten av en förhandsstruktur är; en
modern metaanalys av den frågan dök inte upp.

Vad fältet saknar är en prövning av en försmak med frågor flera veckor
innan undervisningen, i en programmeringskurs, med ett prov som skiljer
mellan de begrepp frågorna gällde och de som de inte gällde.  Visar den
en vinst bara för de tillfrågade begreppen, bör föreläsningens frågor om
det kommande väljas ut efter vad som är viktigast; visar den ingen vinst
alls efter så långt uppehåll, kan försmaken behållas som översikt men
inte motiveras med förtesteffekten.

\needspace{10\baselineskip}
\section{Träffar som rör påståendet}

\Cref{tab:sok-recap-pretest} listar de träffar som rör påståendet.

% --- scholar LaTeX fragment --------------------------------------------
% Generated by: scholar sessions export -f table
% Session: recap-pretest
% \input this file from the body of a document whose
% preamble loads:
%   \usepackage{longtable}
%   \usepackage{booktabs}
%   \usepackage{array}
% ----------------------------------------------------------------------

\begingroup\footnotesize
\begin{longtable}{@{}%
  >{\raggedright\arraybackslash}p{0.44\textwidth}c%
  >{\raggedright\arraybackslash}p{0.13\textwidth}%
  >{\raggedright\arraybackslash}p{0.26\textwidth}@{}}
\caption{Träffar som rör påståendet, förtest, förfrågor och förhandsstrukturer (18 citerade, 56 som stöder påståendet och 83 som kvalificerar eller motsäger det, av 334 unika träffar; de 93 angränsande och 84 felträffarna listas inte). Skälet till varje inkludering står i sista kolumnen; för maskinklassade rader anges modellens konfidens. Databaser: OpenAlex (OA; inte open access), IEEE Xplore (IEEE), Web of Science (WoS), Scopus.}\label{tab:sok-recap-pretest}\\
\toprule
Titel & År & Leverantör & Skäl \\
\midrule
\endfirsthead
\multicolumn{4}{@{}l}{\emph{\tablename~\thetable{} (forts.)}}\\
\toprule
Titel & År & Leverantör & Skäl \\
\midrule
\endhead
\midrule \multicolumn{4}{r@{}}{\emph{forts.\ på nästa sida}}\\
\endfoot
\bottomrule
\endlastfoot
A Meta-Analysis of Advance Organizer Studies & 1983 & OA & \emph{citerad}: förhandsstrukturer \\
A Meta-analysis of the Effects of Advance Organizers on Learning and Retention & 1980 & OA & \emph{citerad}: förhandsstrukturer \\
Beyond the pretesting effect: What happens to the information that is not pretested? & 2019 & OA & \emph{citerad}: bara det tillfrågade \\
Broader benefits of the pretesting effect: Placement matters & 2023 & OA & \emph{citerad}: bara det tillfrågade \\
Can Advance Organizers Influence Meaningful Learning? & 1979 & OA & \emph{citerad}: förhandsstrukturer \\
Does pre-testing promote better retention than post-testing? & 2019 & WoS & \emph{citerad}: före eller efter \\
Guessing as a learning intervention: A meta-analytic review of the prequestion effect & 2024 & Scopus & \emph{citerad}: bara det tillfrågade \\
ManyClasses 2: The Effects of Prequestions on Media Interactions and Learning & 2026 & WoS & \emph{citerad}: i verkliga kurser \\
Prequestioning and Pretesting Effects: a Review of Empirical Research, Theoretical Perspectives, and Implications for Educational Practice & 2023 & Scopus & \emph{citerad}: förtesteffekten \\
Prequestions do not enhance the benefits of retrieval in a STEM classroom & 2017 & OA & \emph{citerad}: i verkliga kurser \\
Pretesting Reduces Mind Wandering and Enhances Learning During Online Lectures & 2020 & WoS, Scopus & \emph{citerad}: förtesteffekten \\
Pretesting Versus Posttesting: Comparing the Pedagogical Benefits of Errorful Generation and Retrieval Practice & 2021 & WoS & \emph{citerad}: före eller efter \\
The Effect of Prequestions on Learning: A Multilevel Meta-Analysis & 2025 & Scopus & \emph{citerad}: bara det tillfrågade \\
The effects of prequestions on classroom learning & 2018 & Scopus & \emph{citerad}: i verkliga kurser \\
The Limited Effects of Prequestions on Learning from Authentic Lecture Videos & 2018 & WoS, Scopus & \emph{citerad}: bara det tillfrågade \\
The pretesting effect comes to full fruition after prolonged retention interval. & 2022 & OA & \emph{citerad}: hur länge \\
The pretesting effect: Do unsuccessful retrieval attempts enhance learning? & 2009 & OA & \emph{citerad}: förtesteffekten \\
The Temporal Dynamics of the Pretesting and Prequestioning Effects & 2026 & WoS & \emph{citerad}: hur länge \\
A Meta-analysis of Effects of the Advance Organizers on Acknowledgment and Retention of Senior Secondary School (SSS) Chemistry & 2011 & OA & stöder påståendet (0.99) \\
Advance organizers in flipped classroom via e-learning management system and the promotion of integrated science process skills & 2019 & OA & stöder påståendet (0.98) \\
Advance organizers: A teaching strategy for achieving significant learning & 2010 & OA & stöder påståendet (0.66) \\
An empirical investigation of the effects on learning and retention of a multiple channel presentation of an advance organizer & 1981 & OA & stöder påståendet (0.78) \\
Application of Augmented Reality technology to promote interactive learning & 2017 & IEEE & stöder påståendet (0.84) \\
Automatic Summarization of Lecture Slides for Enhanced Student PreviewTechnical Report and User Study & 2018 & IEEE & stöder påståendet (0.90) \\
Can Learning Objectives Harness the Power of the Pretesting Effect? (Poster 4) & 2024 & OA & stöder påståendet (0.92) \\
Computer-enriched instruction (CEI) is better for preview material instead of review material: An example of a biostatistics chapter, the central limit theorem & 2010 & Scopus & stöder påståendet (0.98) \\
Decreasing cognitive load for novice EFL learners: Effects of question and descriptive advance organizers in facilitating EFL learners' comprehension of an animation-based content lesson & 2006 & Scopus & stöder påståendet (0.94) \\
Does Covert Prequestioning Enhance Learning? Implications for the Roles of Attention and Retrieval in the Prequestioning Effect & 2026 & Scopus & stöder påståendet (0.99) \\
Duolingo-inspired pretesting with words and pictures improves vocabulary learning & 2026 & Scopus & stöder påståendet (0.99) \\
Effect of Simulation and Pretest Application on Learning in Inhaler Drug Training of Nursing Student: Solomon Experimental Design & 2023 & OA & stöder påståendet (0.97) \\
Effectiveness of virtual simulation-based pre-learning and the mediating role of metacognition in pharmacokinetics laboratory education & 2026 & Scopus & stöder påståendet (0.75) \\
Effects of Advance Organizers on Preschool Children's Learning of Math Concepts & 1978 & OA & stöder påståendet (0.99) \\
Effects of computerized advance organizers on elementary school mathematics learning & 2002 & IEEE & stöder påståendet (0.99) \\
Effects of historical simulations as narrative and graphic advance organizers on Nigerian junior secondary school students' learning outcomes in basic science & 2013 & Scopus & stöder påståendet (0.71) \\
Effects of instructional video playback speed and pre-embedded questions on learning & 2025 & WoS & stöder påståendet (0.96) \\
EFFECTS OF ORIENTING ACTIVITIES AND INSTRUCTIONAL-CONTROL ON LEARNING FACTS AND PROCEDURES FROM INTERACTIVE VIDEO & 1992 & WoS & stöder påståendet (0.91) \\
Enhanced semantic encoding underlies the memory benefit of the pretesting effect & 2026 & OA & stöder påståendet (0.77) \\
Errorful generation and pretesting effects on learning & 2020 & OA & stöder påståendet (0.98) \\
Improving Students' Problem-Solving Skills Using Inquiry Learning Model Combined with Advance Organizer & 2020 & OA & stöder påståendet (0.85) \\
In Defense of Advance Organizers: A Reply to the Critics & 1978 & OA & stöder påståendet (0.82) \\
Instructional Treatments, Advance Organizers and Novice Performance on a Telecommunications Network & 1998 & OA & stöder påståendet (0.94) \\
Interpolated Prequestioning Increases On-Task Thoughts, Dampens Off-Task Thoughts, and Facilitates Learning During Video Lectures & 2026 & WoS & stöder påståendet (0.99) \\
Interpolated pretesting can boost memory of related and distinct prose materials & 2025 & Scopus & stöder påståendet (0.99) \\
Investigating the Benefits of Prequestions on Lecture-Based Learning & 2021 & OA & stöder påståendet (0.92) \\
Metacognitive awareness of the pretesting effect improves with self-regulation support & 2023 & OA & stöder påståendet (0.98) \\
Pan and Rivers (2023) Metacognitive Awareness of the Pretesting Effect Improves with Self-Regulation Support & 2022 & OA & stöder påståendet (0.98) \\
PRACTICAL ANATOMY AS AN ADVANCE ORGANIZER FOR ANATOMY LECTURES: EFFECTIVENESS IN LEARNING FACILITATION FOR DENTAL STUDENTS & 2011 & OA & stöder påståendet (0.89) \\
Prequestions, pretesting, and errorful generation effects on learning outcomes & 2026 & OA & stöder påståendet (0.98) \\
Prereading format and learner proficiency level in L2 reading comprehension & 1990 & Scopus & stöder påståendet (0.91) \\
Pretesting Enhances Learning in the Classroom & 2023 & Scopus & stöder påståendet (1.00) \\
Reinforcing General Chemistry Learning with Structured Daily Recaps: A Pilot Study & 2026 & Scopus & stöder påståendet (0.97) \\
Repeated guessing attempts during acquisition can promote subsequent recall performance. & 2023 & Scopus & stöder påståendet (1.00) \\
Research on the ``Scenario Setup + Virtual Simulation Pre-Study'' Teaching Method for University Physics Experiments Based on Machine Learning - Taking the ``Modification of Ammeter'' Experiment as an Example & 2025 & IEEE & stöder påståendet (0.82) \\
Struggling Prior to a Teaching Event Results in Superior Short-Term Skills Acquisition in Novice Learners & 2020 & Scopus & stöder påståendet (0.99) \\
Teaching With Televised Stories: A Story-Focused Narrative Preview Supports Learning in Young Children & 2020 & Scopus & stöder påståendet (0.99) \\
The effect of a pretest in an interactive, multimodal pretraining system for learning science concepts & 2009 & OA & stöder påståendet (0.98) \\
The effect of inquiry teaching and advance organizers upon student outcomes in science education & 1983 & OA & stöder påståendet (0.95) \\
The Effect of Integrative Prequestions on Learning from Text: An Eye-Tracking Study & 2026 & Scopus & stöder påståendet (0.98) \\
The Effect of Prequestions on Learning from Video Presentations & 2017 & WoS, Scopus & stöder påståendet (1.00) \\
The Effectiveness of Advance Organizers on the Signification of Poetic Images. & 2007 & OA & stöder påståendet (0.96) \\
THE EFFECTIVENESS OF TELEVISION MINIPROGRAMS AS A LEARNING METHOD & 1982 & OA & stöder påståendet (0.66) \\
The Effects of Different Types of Advance Organizers on Classification Learning & 1979 & OA & stöder påståendet (0.99) \\
The effects of prequestions on classroom learning & 2017 & OA & stöder påståendet (0.97) \\
The pretesting effect thrives in the presence of competing information & 2023 & Scopus & stöder påståendet (0.99) \\
The Pretesting Effect: When Tests Improve Memory Even Before Learning & 2026 & WoS & stöder påståendet (1.00) \\
The relationship between segmentation and question location within mobile video platforms for enhancing the ability of recall & 2019 & Scopus & stöder påståendet (0.97) \\
The role of advance organizer on English language learning as a second language & 2010 & OA & stöder påståendet (0.98) \\
The Use of Advance Organizers in the Learning and Retention of Logical Operations and Social Studies Concepts & 1977 & OA & stöder påståendet (0.99) \\
Thinking First Versus Googling First: Preferences and Consequences & 2023 & WoS & stöder påståendet (0.96) \\
Use of an expert concept map as an advance organizer to improve understanding of respiratory failure & 2011 & OA & stöder påståendet (0.96) \\
Using ChatGPT-Generated Prequestions to Improve Memory and Text Comprehension & 2025 & WoS & stöder påståendet (0.99) \\
Video-Mediated Advance Organizers as Facilitators for Greater Listening and Recall & 2010 & OA & stöder påståendet (0.97) \\
Volume, conservation and instruction: A classroom based solomon four group study of conflict & 1981 & OA & stöder påståendet (0.90) \\
Worked-Example Preview for DHH Students in Inclusive Classrooms: A Qualitatively Oriented Multiple-Case Study Design with LASSO-Logistic Survey Modelling & 2026 & Scopus & stöder påståendet (0.79) \\
\textquotedbl{}The Effect of Advance Organizer on Cognitive Load and Learning Achievement\textquotedbl{} & 2006 & OA & kvalificerar eller motsäger påståendet (0.98) \\
``Advance'' advance organizers & 1990 & OA & kvalificerar eller motsäger påståendet (0.99) \\
A meta-analysis of selected advance organizer research reports from 1960-1977 / & 1978 & OA & kvalificerar eller motsäger påståendet (0.94) \\
A Preliminary Investigation of Advance Organizers for a Complex Educational Website & 2005 & OA & kvalificerar eller motsäger påståendet (1.00) \\
Advance Organizers as a Teaching Strategy: A Reply to Barnes and Clawson & 1977 & OA & kvalificerar eller motsäger påståendet (0.96) \\
An integrated model of learning from errors & 2023 & WoS & kvalificerar eller motsäger påståendet (0.91) \\
Attention and the forward testing effect & 2025 & Scopus & kvalificerar eller motsäger påståendet (0.96) \\
Cognitive Aging and Training: The Role of Instructional Coherence and Advance Organizers & 2014 & OA & kvalificerar eller motsäger påståendet (0.93) \\
COGNITIVE STYLE AND ADVANCE ORGANIZERS IN LEARNING AND RETENTION & 1979 & OA & kvalificerar eller motsäger påståendet (0.98) \\
Comparing the effects of different advance organizers on EFL learners' listening comprehension: Key vocabularies, previewing comprehension questions, and multimedia annotations & 2019 & OA & kvalificerar eller motsäger påståendet (0.98) \\
Concurrent Testing Improves Attention Generally and Selectively During Video Lectures & 2026 & WoS & kvalificerar eller motsäger påståendet (0.86) \\
Development and Evaluation of the Kids Count Farm Safety Lesson & 2002 & OA & kvalificerar eller motsäger påståendet (0.90) \\
Differential response to structure of advance organizers in science instruction & 1973 & OA & kvalificerar eller motsäger påståendet (0.99) \\
Differently structured advance organizers lead to different initial schemata and learning outcomes & 2011 & OA & kvalificerar eller motsäger påståendet (0.99) \\
Do curious students learn more science in an immersive virtual reality environment? Exploring the impact of advance organizers and epistemic curiosity & 2022 & OA & kvalificerar eller motsäger påståendet (0.94) \\
Do individual differences in working memory capacity, episodic memory ability, or fluid intelligence moderate the pretesting effect? & 2025 & OA & kvalificerar eller motsäger påståendet (1.00) \\
Do prequestions support learning from text more than learning objectives? & 2026 & Scopus & kvalificerar eller motsäger påståendet (0.95) \\
Does the position and mode of the embedded questions in an instructional video affect learning outcome? A moderation effect of cognitive load level. & 2023 & Scopus & kvalificerar eller motsäger påståendet (0.75) \\
Does viewing a top-down map facilitate spatial learning from navigation in a virtual environment? & 2026 & Scopus & kvalificerar eller motsäger påståendet (0.97) \\
Effect of an Advance Organizer upon Learning for Sixth-Grade Children Maintaining an External Locus of Control Orientation. & 1974 & OA & kvalificerar eller motsäger påståendet (0.81) \\
Effect of Using Photos from Authentic Video as Advance Organizers on Listening Comprehension in an FLES Chinese Class & 2004 & OA & kvalificerar eller motsäger påståendet (0.98) \\
Effects of Advance Organizer Instruction on Preschool Children's Learning of Musical Concepts. & 1992 & OA & kvalificerar eller motsäger påståendet (0.98) \\
Effects of Advance Organizers and Repeated Presentations on Students' Learning & 1996 & OA & kvalificerar eller motsäger påståendet (0.99) \\
Effects of advance organizers on learning and retention from a fully Web-based class & 2007 & OA & kvalificerar eller motsäger påståendet (0.96) \\
Effects Of The Advance Organizer Design For Videos On Mooc Learners & 2020 & OA & kvalificerar eller motsäger påståendet (0.91) \\
Effects of the question wait time and question positions on scientific concept learning from video & 2026 & Scopus & kvalificerar eller motsäger påståendet (0.98) \\
Enhancing Children's Comprehension of a Televised Story Through Previewing & 1990 & OA & kvalificerar eller motsäger påståendet (0.99) \\
Erring on the Side of Caution: Two Failures to Replicate the Derring Effect & 2025 & WoS & kvalificerar eller motsäger påståendet (0.96) \\
Evaluation of a Web-based introductory psychology course: II. Contingency management to increase use of on-line study aids & 2000 & OA & kvalificerar eller motsäger påståendet (0.98) \\
Examining the influence of item exposure and retrieval practice effects on test performance in a large-scale workforce development training programme & 2024 & Scopus & kvalificerar eller motsäger påståendet (0.99) \\
Examining the influence of list composition on the mnemonic benefit of errorful generation & 2025 & Scopus & kvalificerar eller motsäger påståendet (0.99) \\
Feedback is not always better: Interpolating question before learning but no feedback in video lectures facilitating students' performance & 2020 & Scopus & kvalificerar eller motsäger påståendet (0.87) \\
Generative Processing Strengthens Learning Across Errorful and Non-Errorful Contexts & 2026 & Scopus & kvalificerar eller motsäger påståendet (0.99) \\
Guess quality moderates how semantic relatedness influences the pretesting effect. & 2025 & OA & kvalificerar eller motsäger påståendet (0.99) \\
Imagery/Concreteness Attributes of Advance Organizers & 1982 & OA & kvalificerar eller motsäger påståendet (0.98) \\
Implementing virtual reality for emergency training in nephrology: a large-scale study on educational impact and acceptance & 2026 & Scopus & kvalificerar eller motsäger påståendet (0.76) \\
Incorrect Answer in Pretest and Memory Fixation & 2015 & OA & kvalificerar eller motsäger påståendet (0.96) \\
Individual Differences Diminish the Pretest Effect Under Productive Memory Conditions & 2024 & Scopus & kvalificerar eller motsäger påståendet (0.99) \\
Interpolated testing and content pretesting as interventions to reduce task-unrelated thoughts during a video lecture & 2022 & Scopus & kvalificerar eller motsäger påståendet (0.96) \\
Investigations of the Effect of Format of Advance Organizers on Learners' Achievement on Understanding of Science Concepts & 2017 & IEEE & kvalificerar eller motsäger påståendet (0.98) \\
Learning from failure: Errorful generation improves memory for items, not associations & 2019 & Scopus & kvalificerar eller motsäger påståendet (1.00) \\
Let's watch and focus on this: The effects of advance organizers on young children's learning from prosocial and science, technology, engineering, and mathematics (STEM) digital narratives & 2026 & Scopus & kvalificerar eller motsäger påståendet (0.96) \\
Linear and Graphic Advance Organizers: Properties and Processing & 2000 & OA & kvalificerar eller motsäger påståendet (0.95) \\
Making guesses during learning can be beneficial for older adults' memory & 2025 & Scopus & kvalificerar eller motsäger påståendet (0.99) \\
MANIPULATION OF PROCESSING AND MEMORY FOR PROSE THROUGH EXPECTATION AND UNCERTAINTY & 1980 & WoS & kvalificerar eller motsäger påståendet (0.98) \\
On the Role of Generation Rules in Moderating the Beneficial Effects of Errorful Generation & 2021 & WoS & kvalificerar eller motsäger påståendet (1.00) \\
Optimizing the Efficacy of Learning Objectives through Pretests & 2020 & WoS & kvalificerar eller motsäger påståendet (0.98) \\
Prequestions enhance learning, but only when they are remembered. & 2020 & Scopus & kvalificerar eller motsäger påståendet (1.00) \\
Pretesting effects for facts reflect lexical over semantic or structural features & 2026 & Scopus & kvalificerar eller motsäger påståendet (0.99) \\
Pretesting effects on incidental L2 vocabulary learning through reading & 2025 & Scopus & kvalificerar eller motsäger påståendet (0.99) \\
Pretesting Promotes Theme Extraction but Not Generalization From Moral Stories in Young Children & 2026 & WoS & kvalificerar eller motsäger påståendet (1.00) \\
Prior knowledge interacts with the effects of pre-questions and feedback types on learning from videos: Eye-tracking and cognitive load evidence & 2026 & Scopus & kvalificerar eller motsäger påståendet (0.97) \\
QUESTION PREVIEW IN ENGLISH FOR ACADEMIC PURPOSES LISTENING ASSESSMENT: THE EFFECT OF STEM PREVIEW ON DIFFICULTY, ITEM TYPE, AND DISCRIMINATION & 2022 & Scopus & kvalificerar eller motsäger påståendet (0.95) \\
Reading authentic EFL text using visualization and advance organizers in a multimedia learning environment & 2007 & OA & kvalificerar eller motsäger påståendet (0.91) \\
Recall of Advance Organizers as Part of Mathematics Instruction. & 1974 & OA & kvalificerar eller motsäger påståendet (0.96) \\
Research on the Effectiveness of Interactive Learning Activities Based on Rain Classroom & 2020 & IEEE & kvalificerar eller motsäger påståendet (0.89) \\
Scaffolding game-based learning: Impact on learning achievements, perceived learning, and game experiences & 2014 & WoS & kvalificerar eller motsäger påståendet (0.96) \\
Semantic memory and question type : a dive into the pretesting effect & 2026 & OA & kvalificerar eller motsäger påståendet (0.91) \\
Sorry, Am I Intruding? Comparing Performance and Intrusion Rates for Pretested and Posttested Information & 2025 & Scopus & kvalificerar eller motsäger påståendet (0.99) \\
Students' achievement motivation moderates the effects of interpolated pre-questions on attention and learning from video lectures & 2021 & Scopus & kvalificerar eller motsäger påståendet (0.99) \\
Testing before learning: Exploring the robustness of the pretesting effect & 2025 & OA & kvalificerar eller motsäger påståendet (0.98) \\
The benefits of impossible tests: Assessing the role of error-correction in the pretesting effect & 2021 & OA & kvalificerar eller motsäger påståendet (0.99) \\
The difference of pre-class preview time affects online teaching quality & 2020 & IEEE & kvalificerar eller motsäger påståendet (0.88) \\
The Effect of an Advance Organizer, Cognitive Set, and Post Organizer on the Learning and Retention of Written Materials & 1973 & OA & kvalificerar eller motsäger påståendet (0.99) \\
The Effect of Answering Pre-questions and Evaluating Confidence in Preparation on Learning in Classroom Instruction & 2022 & Scopus & kvalificerar eller motsäger påståendet (0.97) \\
The effect of specially constructed advance organizers and postorganizers on mathematics learning & 1978 & OA & kvalificerar eller motsäger påståendet (0.94) \\
The Effects of an Advance Organizer, Text Structure, and a Preview of Structure on the Learning and Retention of Prose Material. & 1980 & OA & kvalificerar eller motsäger påståendet (0.99) \\
The Effects of Using an Advance Organizer on Students' Learning in a Computer Simulation Environment & 1996 & OA & kvalificerar eller motsäger påståendet (0.99) \\
The pretesting effect is robust throughout adulthood, but metacognitive beliefs about pretesting and challenge differ & 2025 & OA & kvalificerar eller motsäger påståendet (0.99) \\
The pretesting effect under divided attention & 2025 & OA & kvalificerar eller motsäger påståendet (0.99) \\
The Pretesting Effect: Exploring the Impact of Feedback and Final Test Timing & 2025 & OA & kvalificerar eller motsäger påståendet (1.00) \\
The relationship of student interest and advance organizer effectiveness & 1988 & OA & kvalificerar eller motsäger påståendet (0.84) \\
The Retention and Usefulness of Concept Maps as Advance Organizers & 2018 & OA & kvalificerar eller motsäger påståendet (0.92) \\
The role of mediators for the pretesting effect & 2024 & WoS & kvalificerar eller motsäger påståendet (0.99) \\
The Trouble With Transfer: Insights From Microgenetic Changes in the Representation of Numerical Magnitude & 2008 & OA & kvalificerar eller motsäger påståendet (0.94) \\
The Value of Using Tests in Education as Tools for Learning-Not Just for Assessment & 2023 & WoS & kvalificerar eller motsäger påståendet (0.91) \\
Type of verbatim question interspersed in text: A new look at the position effect & 1976 & Scopus & kvalificerar eller motsäger påståendet (0.99) \\
Using a digital game as an advance organizer & 2017 & OA & kvalificerar eller motsäger påståendet (0.90) \\
Using Prequestioning as a Hands-On Activity to Support Undergraduate Student Learning & 2025 & Scopus & kvalificerar eller motsäger påståendet (0.99) \\
Using prequestions to enhance learning from reading passages: the roles of question type and structure building ability & 2019 & Scopus & kvalificerar eller motsäger påståendet (1.00) \\
What happens when university students try to answer prequestions that accompany textbook material? & 1990 & Scopus & kvalificerar eller motsäger påståendet (0.93) \\
What Instructional Designers Need To Know about Advance Organizers. & 1998 & OA & kvalificerar eller motsäger påståendet (0.70) \\
When Deliberate Errors Boost Learning: Semantic Constraints on the Derring Effect & 2026 & WoS & kvalificerar eller motsäger påståendet (0.97) \\
\end{longtable}
\endgroup



\chapter{Lär man sig mer av att svara på frågor med telefonen under
  föreläsningen?}
\label{app:recap-polls}
\chapterprecis{Författaren har ännu inte granskat resultaten i den här
  bilagan i sin helhet.}

Lite, och det är omstritt: föreläsningar med frågor som alla besvarar med
ett svarssystem ger i metaanalyserna något bättre resultat än
föreläsningar utan, men vinsten kommer främst av att alla plockar fram
ett svar och av förklaringen efteråt, inte av tekniken.

Frågorna i den här föreläsningen besvaras med telefonen, och alla ser
hur svaren fördelar sig innan svaret visas (den första är
\cref{ovn:recap-start}).  Det är ett svarssystem
(\foreignlanguage{english}{audience response system}), samma idé som de
handhållna knappsatser, \foreignlanguage{english}{clickers}, som
forskningen om sådana system mest har studerat\autocite{Hunsu2016Clickers}.
Frågan den här bilagan besvarar är: \emph{lär sig studenter mer när
föreläsningen har frågor som alla besvarar med ett svarssystem, eller
blir de bara mer engagerade --- och vad avgör om det fungerar?}

\section{Metod}
\label{sec:recap-polls-metod}

Arbetsgången är densamma som i \cref{app:recap-retrieval}, med samma
datum, verktyg och databaser.  Två stödfrågor sökte översikter och
studier i programmering och med kamratlärande
(\foreignlanguage{english}{peer instruction}); två motfrågor sökte
uteblivna effekter, nyhetseffekter och effekter bara på engagemang,
respektive det som avgör om systemet fungerar.  Alla frågor använder
samma sex namn på systemen.  \Cref{tab:recap-polls-queries} visar
frågorna.  Utöver sökningen användes en webbsökning
(\enquote{\foreignlanguage{english}{Hunsu Adesope Bayly 2016
\enquote{meta-analysis of the effects of audience response systems}
pdf}}, 10 oktober 2026) för att hitta en fri kopia av en av
metaanalyserna; den gav i stället en andra, oberoende metaanalys som
inte fanns bland sökningens träffar, och som citeras nedan.

% Author's note (verification and reproduction):
%   scholar session : recap-clickers (S1-S2, M1-M2 on openalex, ieee,
%                     wos, scopus, -l 40; DBLP D1-D4 answered by the
%                     anti-bot page, 0 rows)
%   searched        : 2026-10-10 (s2: key rejected)
%   complement      : WebSearch as quoted in the text, 2026-10-10;
%                     surfaced CastilloManzano2016ARS (not in the session)
%   classification  : as bilaga B
%   provenance      : ltnotes-repetition.bib, keys Hunsu2016Clickers
%                     CastilloManzano2016ARS Chien2016Clickers
%                     Mayer2009Clickers Smith2009PeerDiscussion
%                     Smith2011PeerInstructor ZingaroPorter2014
%                     Akbay2024ARS SerradaSotil2025ARS
%                     CampbellMayer2009Questioning; Deslauriers2019
%                     Boedeker2025 (from bilaga F);
%                     Yang2021Testing, Pan2019Jargon (bilaga B)
\begin{table}[htbp]
  \centering
  \caption{Sökfrågor och antal träffar per databas, 10 oktober 2026,
    verktyget \texttt{scholar}.  S = stödfråga, M = motfråga.  A står
    för (\enquote{audience response system} \textsc{or}
    \enquote{classroom response system} \textsc{or} \enquote{student
    response system} \textsc{or} clicker \textsc{or} clickers
    \textsc{or} \enquote{personal response system}).  Databaser och
    format som i \cref{tab:recap-retrieval-queries}.}
  \label{tab:recap-polls-queries}
  \footnotesize
  \begin{tabular}{@{}lp{0.52\linewidth}rrrr@{}}
    \toprule
    & Nyckelord & OA & IEEE & WoS & Scopus \\
    \midrule
    S1 & A \textsc{and} (learning \textsc{or} achievement) \textsc{and}
      (\enquote{meta-analysis} \textsc{or} \enquote{systematic review}
      \textsc{or} review)
      & 40 & 12 & 40 & 40 \\
    S2 & A \textsc{and} (programming \textsc{or} \enquote{computer
      science} \textsc{or} \enquote{peer instruction})
      & 40 & 40 & 40 & 40 \\
    M1 & A \textsc{and} (\enquote{no significant} \textsc{or} \enquote{no
      effect} \textsc{or} \enquote{novelty effect} \textsc{or}
      engagement) \textsc{and} learning
      & 40 & 40 & 40 & 40 \\
    M2 & A \textsc{and} (\enquote{peer discussion} \textsc{or}
      \enquote{question type} \textsc{or} \enquote{conceptual questions}
      \textsc{or} \enquote{instructor explanation})
      & 40 & 12 & 40 & 40 \\
    \bottomrule
  \end{tabular}
\end{table}

\section{Resultat}
\label{sec:recap-polls-resultat}

\paragraph{Lär sig studenterna mer?}
Lite mer, enligt två metaanalyser.  Den ena konstaterar att tidigare
resultat var \enquote{largely mixed and inconclusive} och finner
\enquote{small but significant effects} på både kunskap och inställning,
beroende av ämne, gruppstorlek och hur frågorna används
\parencite{Hunsu2016Clickers}.
Den andra, med 51 studier och nästan 15\,000 deltagare, fann en effekt
på provresultat som var mindre på universitet, \enquote{g = 0.22}, än i
skolan, \enquote{g = 0.48}\autocite[109]{CastilloManzano2016ARS}.  Den nyaste,
med 68 studier, fann en liten effekt på kunskap, \enquote{g = 0.47},
större i skolan än på högskolan\autocite{Akbay2024ARS}.  En
systematisk genomgång av studier i högre utbildning --- där verktyg av
samma slag som det i den här föreläsningen dominerar --- fann tvärtom
metoderna svaga och \enquote{no consensus in the literature regarding
how ARS contribute to academic performance}\autocite{SerradaSotil2025ARS}.
En äldre översikt menar att fördelen jämfört med vanliga föreläsningar
\enquote{stands on firm empirical ground}\autocite{Chien2016Clickers}.

\paragraph{Är det systemet eller frågorna?}
Främst frågorna och vad som görs med dem.  Klassrumsmetaanalysen i
\cref{app:recap-retrieval} fann att testeffekten inte beror på om svaret
ges med penna, på nätet eller med ett svarssystem\autocite{Yang2021Testing}
--- vinsten kommer av att plocka fram svaret.  I två
laboratorieförsök med en föreläsning gav fyra frågor med svarssystem,
klassens svarsfördelning och lärarens förklaring bättre resultat än
samma innehåll framfört som påståenden med samma förklaring, på ett
minnesprov (\enquote{d = 1.23}) och på ett överföringsprov
(\enquote{d = 0.74}), men inte på de övriga
proven\autocite{CampbellMayer2009Questioning}.
Ett försök i en stor
föreläsningskurs fann ändå att samma frågor gav bättre provresultat när
alla svarade med knappsats (\enquote{d = 0.38}) än när de ställdes utan
den, och att frågor utan knappsats inte gav något jämfört med inga
frågor alls \parencite{Mayer2009Clickers}; försöket jämförde tre olika år
av samma kurs, inte slumpade grupper.  Författarna förklarar vinsten
med att studenterna med knappsats \enquote{are more cognitively engaged
during learning} \parencite{Mayer2009Clickers}.

\paragraph{Vad avgör om det fungerar?}
Att svaret följs av förklaring och resonemang.  Att låta studenterna
förklara och motivera sina svar är det upplägg som hänger ihop med de
största effekterna \parencite{Chien2016Clickers}.  Diskussion med grannen
efter en omröstning ökar förståelsen, även när ingen i gruppen först
visste svaret \parencite{Smith2009PeerDiscussion}, och diskussion följd av
lärarens förklaring ger mer än vardera för
sig\autocite{Smith2011PeerInstructor}.
I programmeringskurser med kamratlärande gav lärarens genomgång efter
diskussionen en tydlig extra vinst, störst för svagare studenter och
svårare frågor\autocite{ZingaroPorter2014}.

\paragraph{Lär man sig mindre än man tror?}
Inte nödvändigtvis, men det kan kännas så.  I ett lottat försök i en
fysikkurs lärde sig studenterna i aktivt undervisade lektioner mer än i
flytande föreläsningar, men upplevde att de lärt sig mindre och föredrog
föreläsningarna\autocite[19251]{Deslauriers2019}; i ett liknande försök
med läkarstudenter upplevde de aktiva studenterna tvärtom att de lärt sig
mer\autocite{Boedeker2025}.  Studenternas omdöme om omröstningarna säger
alltså inte säkert något om vad de lär sig.%
% Bilaga F is in the teacher's edition only.
\ifbool{ltnote}{  Båda försöken redovisas i
\cref{app:recap-flippat}.}{}

\section{Diskussion och begränsningar}
\label{sec:recap-polls-diskussion}

Effekten av själva systemet är liten, och den mesta vinsten går att
förklara med testeffekten i \cref{app:recap-retrieval} och med det som
görs efter omröstningen.  Ingen av källorna har studerat omröstningar
via telefon i en programmeringsföreläsning för nybörjare, och studien
i programmeringskurser mätte vinsten inom en fråga, inte på ett prov.
Tre av metaanalyserna och fyra av studierna är lästa bara i
sammanfattningen, och den första metaanalysens effektstorlekar står inte
i sammanfattningen och anges därför inte.  Metaanalyserna är inte
oberoende av varandra: de bygger delvis på samma studier.

\section{Slutsats}
\label{sec:recap-polls-slutsats}

Ja, men lite och omstritt: föreläsningar med frågor som alla besvarar
med ett svarssystem ger i metaanalyserna något bättre resultat än
föreläsningar utan, mindre på högskolan än i skolan, medan en
systematisk genomgång av högskolestudierna finner dem för svaga för en
slutsats.  Vinsten kommer främst av
att alla plockar fram ett svar och av förklaringen som följer, inte av
tekniken.  Föreläsningen kan därför säga att omröstningen får alla att
svara och visar vad gruppen tror, och att det är genomgången efter den
som ska ge förståelsen.

\section{Fortsatt arbete}
\label{sec:recap-polls-fortsatt}

Semantic Scholar och DBLP saknas, och DBLP indexerar de konferenser där
kamratlärande i programmering mest har publicerats; sökningen S2 bör
göras om där.  Sju källor är lästa bara i sammanfattningen.  Den andra
metaanalysen kom från en webbsökning, inte från sökningen, vilket
visar att sökningens 40 träffar per fråga inte räckte för att fånga alla
översikter.

Vad fältet saknar är ett försök med omröstningar via telefon i en
programmeringsföreläsning, där samma frågor ställs med och utan
omröstning, med och utan diskussion efteråt, och med ett prov som mäter
förståelse av programkod.  Visar det att omröstningen inte tillför något
utöver frågan, kan föreläsningen lika gärna ställa frågorna muntligt;
visar det att diskussionen efteråt är det som ger vinsten, bör
föreläsningen ge tid åt den.

\needspace{10\baselineskip}
\section{Träffar som rör påståendet}

\Cref{tab:sok-recap-clickers} listar de träffar som rör påståendet.

% --- scholar LaTeX fragment --------------------------------------------
% Generated by: scholar sessions export -f table
% Session: recap-clickers
% \input this file from the body of a document whose
% preamble loads:
%   \usepackage{longtable}
%   \usepackage{booktabs}
%   \usepackage{array}
% ----------------------------------------------------------------------

\begingroup\footnotesize
\begin{longtable}{@{}%
  >{\raggedright\arraybackslash}p{0.44\textwidth}c%
  >{\raggedright\arraybackslash}p{0.13\textwidth}%
  >{\raggedright\arraybackslash}p{0.26\textwidth}@{}}
\caption{Träffar som rör påståendet, svarssystem i föreläsningar (9 citerade, 131 som stöder påståendet och 81 som kvalificerar eller motsäger det, av 360 unika träffar; de 117 angränsande och 22 felträffarna listas inte). Skälet till varje inkludering står i sista kolumnen; för maskinklassade rader anges modellens konfidens. Databaser: OpenAlex (OA; inte open access), IEEE Xplore (IEEE), Web of Science (WoS), Scopus.}\label{tab:sok-recap-clickers}\\
\toprule
Titel & År & Leverantör & Skäl \\
\midrule
\endfirsthead
\multicolumn{4}{@{}l}{\emph{\tablename~\thetable{} (forts.)}}\\
\toprule
Titel & År & Leverantör & Skäl \\
\midrule
\endhead
\midrule \multicolumn{4}{r@{}}{\emph{forts.\ på nästa sida}}\\
\endfoot
\bottomrule
\endlastfoot
A meta-analysis of the effects of audience response systems (clicker-based technologies) on cognition and affect & 2015 & OA & \emph{citerad}: metaanalyser och översikter \\
Clickers in college classrooms: Fostering learning with questioning methods in large lecture classes & 2008 & OA & \emph{citerad}: frågorna eller systemet \\
Combining Peer Discussion with Instructor Explanation Increases Student Learning from In-Class Concept Questions & 2011 & WoS & \emph{citerad}: diskussion och förklaring \\
Do audience response systems truly enhance learning and motivation in higher education? A systematic review & 2025 & Scopus & \emph{citerad}: metaanalyser och översikter \\
Do we click in the right way? A meta-analytic review of clicker-integrated instruction & 2015 & OA & \emph{citerad}: metaanalyser och översikter \\
Peer Instruction in computing: The value of instructor intervention & 2014 & WoS & \emph{citerad}: diskussion och förklaring \\
Questioning as an instructional method: Does it affect learning from lectures? & 2008 & OA & \emph{citerad}: frågorna eller systemet \\
Re-Examining the Effect of Audience Response Systems on Learning Outcomes: Evidence from the Last Decade & 2024 & Scopus & \emph{citerad}: metaanalyser och översikter \\
Why Peer Discussion Improves Student Performance on In-Class Concept Questions & 2009 & OA & \emph{citerad}: diskussion och förklaring \\
<CLICK HERE> FOR INTERACTION: ABOUT IMPLEMENTING CLICKERS IN THE CLASSROOM - EXPERIENCES AND OPINIONS & 2011 & WoS & stöder påståendet (0.77) \\
A Collaborative and Competition-based Learning Approach for Enhanced Student Engagement and Performance; Pre, During, and Post-Covid Perspectives & 2024 & WoS & stöder påståendet (0.96) \\
A Comparison And Evaluation Of Personal Response Systems In Introductory Computer Programming & 2020 & OA & stöder påståendet (0.67) \\
A Hybrid Student-centered Learning Instructional Approaches for a Lecture Based Engineering Class & 2017 & IEEE & stöder påståendet (0.92) \\
A New Approach to Perform Individual Assessments at Higher Education Using Gamification Systems & 2023 & Scopus & stöder påståendet (0.94) \\
A Review of Personal Response Systems in Higher Education & 2022 & OA & stöder påståendet (0.98) \\
A systematic review of classroom technologies supporting student centered teaching & 2025 & Scopus & stöder påståendet (0.91) \\
Active learning in environmental engineering: Combining interactive platforms and project-based approaches to boost engagement and academic performance & 2025 & Scopus & stöder påståendet (0.90) \\
Active learning strategies and student engagement in undergraduate medical education: A systematic review (2015--2025) & 2026 & Scopus & stöder påståendet (0.84) \\
Addition of an audience response system in a biomedical science course improves learning environment and student performance & 2024 & WoS & stöder påståendet (0.99) \\
Addressing common errors and misconceptions in integral calculus with clickers and classroom voting & 2020 & IEEE & stöder påståendet (0.93) \\
Advancing chemical engineering education: Amplifying active learning with Wooclap's innovative pedagogical techniques & 2025 & Scopus & stöder påståendet (0.97) \\
AN EVALUATION OF STUDENT RESPONSE SYSTEMS FROM THE VIEWPOINT OF INSTRUCTORS AND STUDENTS & 2011 & OA & stöder påståendet (0.95) \\
An Evidence Based Learning and Teaching Strategy for Computer Science Classrooms and Its Extension into a Mobile Classroom Response System & 2014 & IEEE & stöder påståendet (0.89) \\
Asking the classroom with technology: a systematic literature review & 2025 & WoS & stöder påståendet (0.97) \\
Assessing Information Literacy Comprehension in First-Year Students & 2011 & OA & stöder påståendet (0.93) \\
Assessing Multimedia Influences on Student Responses Using a Personal Response System & 2011 & OA & stöder påståendet (0.98) \\
Audience response systems as an interim measure of quality using Bloom's taxonomy of learning outcomes & 2022 & Scopus & stöder påståendet (0.88) \\
Augmented reality meets Peer instruction & 2024 & Scopus & stöder påståendet (0.98) \\
Augmenting Experiential Learning with a Game-based Student Response System & 2018 & IEEE & stöder påståendet (0.92) \\
Behind the scenes: Unpacking student discussion and critical reflection in lectures & 2020 & Scopus & stöder påståendet (0.67) \\
BOARD \# 41: Enhancing Student Engagement and Learning Outcomes: Comparing Interactive Simulations with Traditional Clicker Questions in Introductory Engineering Courses & 2025 & Scopus & stöder påståendet (0.79) \\
Can Technology Improve Large Class Learning? The Case of an Upper-Division Business Core Class & 2013 & OA & stöder påståendet (0.99) \\
Case study: Applying a motivation technique to improve student performance in higher education & 2023 & Scopus & stöder påståendet (0.94) \\
Catalyzing Learner Engagement using Cutting-Edge Classroom Response Systems in Higher Education & 2025 & Scopus & stöder påståendet (0.91) \\
Clicker-integrated instruction and conventional instruction: The comparative evaluations of students' performances in chemistry & 2025 & Scopus & stöder påståendet (0.91) \\
Clickers in the Biology Classroom: Strategies for Writing and Effectively Implementing Clicker Questions That Maximize Student Learning & 2020 & Scopus & stöder påståendet (0.83) \\
Clickers: A New Teaching Aid with Exceptional Promise & 2006 & OA & stöder påståendet (0.98) \\
Clicking Our Way To Class Discussion & 2010 & OA & stöder påståendet (0.97) \\
Collaborative learning in engineering physics tutorials & 2018 & Scopus & stöder påståendet (0.79) \\
Comparative effect of interactive mobiles (clickers) and communicative approach on the learning outcomes of the educationally disadvantaged Nigerian pupils in ESL classrooms & 2013 & OA & stöder påståendet (0.95) \\
Comparing the effectiveness of Jigsaw and Audience Response System on patient safety competency in nursing students: A quasi-experimental study & 2025 & Scopus & stöder påståendet (0.98) \\
Comparison of Student Responses to Easy and Difficult Thermodynamics Conceptual Questions during Peer Instruction & 2011 & WoS & stöder påståendet (0.97) \\
Correlation between Academic Achievement Results and Students' Perceptions in Instant Response System-Based Language Learning Classes at the University & 2024 & WoS & stöder påståendet (0.91) \\
Developing design principles and task types for classroom response system tasks in mathematics & 2022 & Scopus & stöder påståendet (0.91) \\
Development, validation and online and in-person implementation of clicker question sequence on change of basis & 2023 & WoS & stöder påståendet (0.92) \\
Do feedback strategies improve students' learning gain?-Results of a randomized experiment using polling technology in physics classrooms & 2021 & WoS & stöder påståendet (0.99) \\
Do Personal Response Systems Improve Learning? & 2016 & OA & stöder påståendet (0.88) \\
Effective learning in science: The use of personal response systems with a wide range of audiences & 2010 & OA & stöder påståendet (0.95) \\
Effects of gamification on the engagement and translation proficiency of the Chinese EFL learners: a case of Quizizz & 2026 & WoS & stöder påståendet (0.97) \\
Effects of Technology-Enhanced Learning on Student Engagement in Basic Medical Education: A Systematic Review & 2025 & Scopus & stöder påståendet (0.94) \\
Engaging students inside the classroom to increase learning & 2015 & IEEE & stöder påståendet (0.85) \\
Enhanced Conceptual Learning with Real Time Student-Generated Data and Visualization & 2025 & WoS & stöder påståendet (0.90) \\
Enhancing Civil Engineering Education: A Comprehensive Analysis of Student Perspectives on Technology-Integrated Learning & 2024 & Scopus & stöder påståendet (0.96) \\
Enhancing Grammatical Accuracy in EFL Students: A Game-Based Approach Utilizing Peer Response Displayed through an Interactive Application & 2024 & IEEE & stöder påståendet (0.85) \\
Enhancing Learning in the Classroom -- One Click @ a Time & 2011 & OA & stöder påståendet (0.93) \\
Enhancing mandarin learning through classroom response apps: A comparison of cognitive and affective factors in touch and shaking responses & 2024 & WoS & stöder påståendet (0.92) \\
Enhancing Second Language Skills Development Using Students Response System & 2012 & WoS & stöder påståendet (0.98) \\
Evaluating classroom response systems in engineering education: Which metrics better reflect student performance? & 2025 & Scopus & stöder påståendet (0.97) \\
Evaluating effectiveness of active learning in computer science using metacognition & 2014 & IEEE & stöder påståendet (0.97) \\
Experiences with the use of Classroom Response System (CRS) in a Chemistry Foundation Course & 2010 & OA & stöder påståendet (0.62) \\
EXPLORATION OF THE USE OF HANDHELD PERSONAL RESPONSE SYSTEMS WITH FIRST YEAR ACCOUNTANCY STUDENTS FOR DEEP LEARNING AND UNDERSTANDING & 2012 & OA & stöder påståendet (0.86) \\
Exploring Learning-Oriented Assessment in Enhancing Students' Lexical Fluency Through MALL & 2025 & Scopus & stöder påståendet (0.98) \\
Exploring the Efficacy of Student Response System in a Sub-Saharan African Country: A Sociocultural Perspective & 2012 & WoS & stöder påståendet (0.97) \\
Faculty and student experiences with clickers: a qualitative exploration of engaging students in higher-level thinking & 2008 & OA & stöder påståendet (0.96) \\
Formative teaching: A Conversational Framework for evaluating the impact of Response Technology on student experience, engagement and achievement & 2009 & IEEE & stöder påståendet (0.99) \\
High Enrollment, Early Engineering Courses And The Personal Response System & 2020 & OA & stöder påståendet (0.88) \\
HOW DOES A CLASSROOM INTERACTION SYSTEM AFFECT STUDENT & 2007 & OA & stöder påståendet (0.96) \\
How quiz-based tools can improve students' engagement and participation in the classroom & 2014 & IEEE & stöder påståendet (0.98) \\
How using a response system in blended synchronous seminars encourages online and onsite student participation & 2024 & WoS & stöder påståendet (0.97) \\
iCAP: A Classroom Engagement Tool for Introductory Programming Courses & 2023 & WoS & stöder påståendet (0.88) \\
Identifying Clicker Questions that Provoke Rich Discussions in Introductory Statistics & 2022 & Scopus & stöder påståendet (0.96) \\
Identifying Key Features of Effective Active Learning: The Effects of Writing and Peer Discussion & 2014 & WoS & stöder påståendet (0.96) \\
Impact of a game-based tool on student engagement in a foreign language course: a three-term analysis & 2024 & WoS & stöder påståendet (0.94) \\
IMPACT OF IMPLEMENTING CLASS RESPONSE SYSTEM IN ELECTRONICS ENGINEERING COURSES TOWARDS STUDENTS' ENGAGEMENT IN CLASS & 2013 & WoS & stöder påståendet (0.92) \\
Implementation and evaluation of different types of peer learning instruction in a MATLAB programming course & 2016 & Scopus & stöder påståendet (0.68) \\
Implications of using classroom response systems (CRS) on learning performance: An experience of learning analytics & 2022 & Scopus & stöder påståendet (0.97) \\
Improving learning and engagement for students in large classes & 2007 & IEEE & stöder påståendet (0.82) \\
Improving student engagement in higher education through mobile-based interactive teaching model using socrative & 2017 & IEEE & stöder påståendet (0.94) \\
Incorporating Immediate Online Feedback System and Concept Map into Electrical Circuits Course & 2022 & WoS & stöder påståendet (0.94) \\
Incorporating Students' Self-Efficacy and Subject Value in the Evaluation of Audience Response Systems & 2015 & IEEE & stöder påståendet (0.95) \\
Integrating Quiz-Based Self-Assessment Into a Microbiology Practical Course: Effects on Student Perception and Learning Outcomes & 2026 & WoS & stöder påståendet (0.95) \\
Integrating Student Response Technology into a Large Undergraduate Course: Students' Perceptions of their Motivations and Feedback & 2024 & WoS & stöder påståendet (0.93) \\
Interactive question based learning methodology and clickers: Fundamentals of computer science course case study & 2013 & IEEE & stöder påståendet (0.99) \\
Interactive surveys during online lectures for IT students & 2023 & Scopus & stöder påståendet (0.65) \\
Interactive voting systems as promoters of the transformation of individual and collective activity in master classes in higher education & 2024 & Scopus & stöder påståendet (0.77) \\
Investigating the effects of peer instruction on preservice mathematics teachers' achievements in statistics and probability & 2018 & WoS & stöder påståendet (0.95) \\
Investigation into the Impact of Audience Response Devices on Short- and Long-term Content Retention & 2013 & WoS & stöder påståendet (0.99) \\
Learners' Self-Regulation in an Interactive Response System-Aided Flipped Classroom & 2017 & IEEE & stöder påståendet (0.93) \\
Leveraging Intelligent Systems to Enhance Student Engagement and Attendance in the Classroom & 2026 & Scopus & stöder påståendet (0.95) \\
Modeling Active Engagement Pedagogy through Classroom Response Systems in a Physics Teacher Education Course & 2013 & OA & stöder påståendet (0.84) \\
Multi-institutional assessment of Peer Instruction implementation and impacts using the Framework for Interactive Learning in Lectures & 2025 & Scopus & stöder påståendet (0.92) \\
Online Tests Based on Contributions Provided by Teachers and Students during Face to Face Lectures & 2015 & IEEE & stöder påståendet (0.94) \\
Peer Assessment of Oral Presentation: An Investigative Study of Using Clickers in First-Year Civil Engineering Class of a Reputed Engineering Institution & 2016 & IEEE & stöder påståendet (0.93) \\
Peer assessment of oral presentations using clickers: the student experience & 2009 & OA & stöder påståendet (0.94) \\
Peer discussions and response technology: Short interventions, considerable gains & 2017 & Scopus & stöder påståendet (0.99) \\
Peer Instruction in Online Synchronous Software Engineering - Findings from fine-grained clicker data & 2021 & Scopus & stöder påståendet (0.97) \\
Peer Instruction: Using Isomorphic Questions to Document Learning Gains in a Small Statics Class & 2016 & Scopus & stöder påståendet (0.99) \\
Perceptions of First-year MBBS Students towards Kahoot! as a Formative Assessment Tool in Medical Education: A Quasi-experimental Study from Howrah, India & 2026 & WoS & stöder påståendet (0.94) \\
Personal response systems and learning performance: The mediating role of learners' engagement & 2019 & OA & stöder påståendet (0.99) \\
Personal response systems through the prism of students' experiences & 2020 & OA & stöder påståendet (0.94) \\
Preliminary Results Of Using Personal Response Systems (Clickers) In A Conceptual Physics Course & 2020 & OA & stöder påståendet (0.98) \\
Promoting Intercultural Understanding and Motivation in Language Learning Through Kahoot! & 2025 & Scopus & stöder påståendet (0.90) \\
Promoting student-centered active learning in lectures with a personal response system & 2009 & OA & stöder påståendet (0.99) \\
Provocation to Learn - A Study in the Use of Personal Response Systems in Information Literacy Instruction & 2008 & OA & stöder påståendet (0.78) \\
Research-Based Implementation of Peer Instruction: A Literature Review & 2015 & WoS & stöder påståendet (0.99) \\
Sequential Exercises and Personal Response System in Project Management Courses & 2016 & OA & stöder påståendet (0.93) \\
Small Group Test of the Personal Response System (PRS) in a Behavioral Science Graduate Research Methods Course & 2006 & OA & stöder påståendet (0.95) \\
Student Opinions and Preferences Regarding Personal Response Systems in the Graduate Physical Therapy Classroom: A Mixed-Methods Inquiry & 2012 & OA & stöder påståendet (0.95) \\
Student perceptions of e-learning components within a Masters level accounting module & 2012 & IEEE & stöder påståendet (0.89) \\
Student Response System in Dental Students' Education. Using a Student Response System and Peer Discussion to Raise the Awareness of the Importance of Good Professional Communication Skills in Practice Periods & 2020 & Scopus & stöder påståendet (0.87) \\
Student Use of Classroom Response Systems to Promote Active Learning & 2019 & Scopus & stöder påståendet (0.89) \\
Student's Perceptions of Technology-Mediated Open-Text Questions in the Classroom: A Case Study & 2025 & Scopus & stöder påståendet (0.94) \\
Student-active lectures in legal education & 2024 & WoS & stöder påståendet (0.95) \\
Study on the Learners' Performance Evaluation in Online Teaching with Personal Response Systems & 2021 & IEEE & stöder påståendet (0.98) \\
Synchronous and Asynchronous Formative E-Assessments for Scalable Mentoring in Engineering Education & 2023 & IEEE & stöder påståendet (0.72) \\
Tailoring Clicker Technology to Problem-Based Learning: What's the Best Approach? & 2017 & Scopus & stöder påståendet (0.94) \\
Teaching SOLID Software Design Principles Using Peer Instruction-A Pilot Study & 2024 & Scopus & stöder påståendet (0.78) \\
The Benefits of Using Clickers in Small-Enrollment Seminar-Style Biology Courses & 2011 & WoS & stöder påståendet (0.98) \\
The effect of using personal response system on 6th grade students' achievement and attitudes towards science and technology & 2022 & OA & stöder påståendet (0.96) \\
The effective use of clickers in freshmen classrooms & 2011 & IEEE & stöder påståendet (0.99) \\
The Effectiveness of Library Instruction: Do Student Response Systems (Clickers) Enhance Learning? & 2010 & OA & stöder påståendet (0.99) \\
THE EFFECTIVENESS OF PEER INSTRUCTION USING PERSONAL RESPONSE SYSTEMS ON THE ACADEMIC PERFORMANCE AND ATTITUDES OF STUDENTS IN AN INTRODUCTORY PSYCHOLOGY CLASS & 2009 & OA & stöder påståendet (0.68) \\
The impact of individual and shared clicker use on students' collaborative learning & 2015 & OA & stöder påståendet (0.99) \\
The Impact of Student Response Systems (SRS) on Student Achievements: A University-Scale Study with Deep Exploratory Data Analysis (EDA) & 2023 & Scopus & stöder påståendet (0.98) \\
The Impact of Technology-Assisted \textquotedbl{}scaffolding\textquotedbl{} on Student Learning in General Chemistry & 2019 & Scopus & stöder påståendet (0.82) \\
The Integration of Personal Response Systems into Online Flipped Classroom: Student Performance and Perception & 2021 & IEEE & stöder påståendet (0.99) \\
Translating Neuroscience: When is the use of Clickers Effective for Student Learning? & 2013 & OA & stöder påståendet (0.94) \\
Use of automated response systems in the small sized class & 2013 & OA & stöder påståendet (0.70) \\
Use of Online In-Class Quizzes to Motivate and Engage Students in Courses on Semiconductor Devices and Integrated Circuits & 2022 & IEEE & stöder påståendet (0.95) \\
Use of remote response devices: an effective interactive method in the long- term learning & 2014 & OA & stöder påståendet (1.00) \\
Using a personal response system as an in-class assessment tool in the teaching of basic college chemistry & 2013 & OA & stöder påståendet (0.95) \\
Using clickers in lectures to help identify and teach the control topics students find difficult & 2014 & IEEE & stöder påståendet (0.90) \\
Using Clickers to Address Common Errors and Misconceptions in Statistics & 2023 & Scopus & stöder påståendet (0.94) \\
Using Clickers to Facilitate Development of Problem-Solving Skills & 2011 & OA & stöder påståendet (0.99) \\
Using Mentimeter to Enhance Students' Engagement in Large University Classes - A Survey of Students' Perceptions & 2023 & IEEE & stöder påståendet (0.98) \\
Using Peer Discussion Facilitated by Clicker Questions in an Informal Education Setting: Enhancing Farmer Learning of Science & 2012 & WoS & stöder påståendet (0.96) \\
Using technology to enhance the implementation of peer discussion in science education & 2013 & OA & stöder påståendet (0.98) \\
Virtual Tools and Student Engagement: Insights from Engineering Education & 2025 & Scopus & stöder påståendet (0.96) \\
What are Middle School Students Talking About During Clicker Questions? Characterizing Small-Group Conversations Mediated by Classroom Response Systems & 2016 & WoS & stöder påståendet (0.98) \\
Å lære mekanikk ved bruk av et elektronisk \textquotedbl{}Personal Response System\textquotedbl{} & 2008 & OA & stöder påståendet (0.75) \\
``Clickers'' as Catalysts for Transformation of Teachers & 2010 & OA & kvalificerar eller motsäger påståendet (0.96) \\
A campus-wide investigation of clicker implementation: The status of peer discussion in STEM classes & 2016 & Scopus & kvalificerar eller motsäger påståendet (0.92) \\
A literature review on the influence of Kahoot! On learning outcomes, interaction, and collaboration & 2021 & WoS & kvalificerar eller motsäger påståendet (0.99) \\
A review of literature reports of clickers applicable to college chemistry classrooms & 2008 & OA & kvalificerar eller motsäger påståendet (0.99) \\
A study on the influence of rich versus traditional classroom response system (CRS) questions on concept retention & 2012 & IEEE & kvalificerar eller motsäger påståendet (0.87) \\
A systematic literature review of web-based student response systems: Advantages and challenges & 2022 & Scopus & kvalificerar eller motsäger påståendet (0.99) \\
Advantages and disadvantages of clicker use in education & 2017 & Scopus & kvalificerar eller motsäger påståendet (0.97) \\
All Students Can Benefit from a Personal Response System Pedagogy that Encourages Active Engagement Yet Lowers Barriers to Implementation & 2025 & Scopus & kvalificerar eller motsäger påståendet (0.94) \\
An exploration of teachers' perceptions towards student engagement in technology-assisted learning classrooms at a sino-foreign university in China & 2026 & Scopus & kvalificerar eller motsäger påståendet (0.87) \\
Analysis of participation rates in Poll Everywhere questions and academic performance in a veterinary biochemistry and metabolism course & 2024 & WoS & kvalificerar eller motsäger påståendet (0.96) \\
Assessing the Use of Kahoot! in an Undergraduate General Chemistry Classroom & 2022 & Scopus & kvalificerar eller motsäger påståendet (0.94) \\
Audience Response System (ARS) Use in the SCORE (Surgical Council on Resident Education) Surgery Training Curriculum: A Mixed Methodology Study & 2023 & WoS & kvalificerar eller motsäger påståendet (0.93) \\
Audience response system utilization outcome in secondary school environment & 2014 & IEEE & kvalificerar eller motsäger påståendet (0.99) \\
Automated Intelligent Feedback Based Learning in a Software Development Project Management Course & 2023 & WoS, Scopus & kvalificerar eller motsäger påståendet (0.98) \\
Balancing Engagement and Integrity: Rethinking Assessment Strategies in Team-Based Learning & 2026 & Scopus & kvalificerar eller motsäger påståendet (0.99) \\
Benefits of Electronic Audience Response Systems on Student Participation, Learning, and Emotion & 2007 & OA & kvalificerar eller motsäger påståendet (0.97) \\
Can a \$30 Piece of Plastic Improve Learning? An Evaluation of Personal Response Systems in Large Classroom Settings & 2004 & OA & kvalificerar eller motsäger påståendet (0.96) \\
Challenges in addressing student difficulties with quantum measurement of two-state quantum systems using a multiple-choice question sequence in online and in-person classes & 2023 & WoS & kvalificerar eller motsäger påståendet (0.92) \\
Clicker interventions at university lectures and the feedback gap & 2019 & Scopus & kvalificerar eller motsäger påståendet (0.99) \\
Clickers and formative feedback at university lectures & 2017 & Scopus & kvalificerar eller motsäger påståendet (0.91) \\
Clickers as Data Gathering Tools and Students' Attitudes, Motivations, and Beliefs on Their Use in this Application & 2009 & OA & kvalificerar eller motsäger påståendet (0.95) \\
Clickers in the Large Classroom: Current Research and Best-Practice Tips & 2007 & OA & kvalificerar eller motsäger påståendet (0.97) \\
Clicking with your audience: Evaluating the use of personal response systems in library instruction & 2011 & OA & kvalificerar eller motsäger påståendet (0.99) \\
Cloze tests as retrieval practice activities: evaluating their integration with audience response systems in K-12 schools & 2025 & Scopus & kvalificerar eller motsäger påståendet (0.99) \\
Conceptual question response times in Peer Instruction classrooms & 2014 & OA & kvalificerar eller motsäger påståendet (0.86) \\
Does Displaying the Class Results Affect Student Discussion during Peer Instruction? & 2010 & OA & kvalificerar eller motsäger påståendet (0.99) \\
Does Peer Instruction Enhance Impact of Audience Response Systems for Content Retention? & 2025 & Scopus & kvalificerar eller motsäger påståendet (0.99) \\
Does Question Complexity Moderate the Effects of Clicker Use on Student Learning? & 2023 & Scopus & kvalificerar eller motsäger påståendet (0.97) \\
Effect of personal response systems on student perception and academic performance in courses in a health sciences curriculum & 2011 & WoS & kvalificerar eller motsäger påståendet (0.98) \\
Effectiveness of real-time classroom interactive competition on academic performance: a systematic review and meta-analysis & 2023 & Scopus & kvalificerar eller motsäger påståendet (0.99) \\
Effectiveness Of Using Personal Response Systems In A Conceptual Physics Course & 2020 & OA & kvalificerar eller motsäger påståendet (0.95) \\
Effects of Student Response Systems on Learning Outcomes & 2026 & WoS & kvalificerar eller motsäger påståendet (0.98) \\
Electronic voting systems for lectures then and now: A comparison of research and practice & 2007 & OA & kvalificerar eller motsäger påståendet (0.97) \\
Enhancing Check-Reinforce Introduction With a Class Response System: The C2RI Method - A Five-Year Study & 2020 & IEEE & kvalificerar eller motsäger påståendet (0.91) \\
Enhancing Synchronous Collaborative Learning with AI-Supported Audience Response Systems: The EchoQuiz Approach & 2026 & Scopus & kvalificerar eller motsäger påståendet (0.97) \\
Experience with a multiple-choice audience response system in an engineering classroom & 2005 & IEEE & kvalificerar eller motsäger påståendet (0.94) \\
Exploring Peer Instruction: Should Cohort Clicker Responses Appear during or after Polling? & 2019 & Scopus & kvalificerar eller motsäger påståendet (0.98) \\
Facilitating Active Collaborative Learning in Medical Education; a Literature Review of Peer Instruction Method & 2023 & Scopus & kvalificerar eller motsäger påståendet (0.97) \\
Formative assessment in statistics education: A scoping review & 2026 & Scopus & kvalificerar eller motsäger påståendet (0.76) \\
From emotional safety to agentic engagement: implementing Mentimeter in undergraduate TESOL courses & 2026 & Scopus & kvalificerar eller motsäger påståendet (0.98) \\
Gamification and Active Learning in Agricultural Engineering: Evaluating the Impact of an Interactive Response System on Academic Performance and Stress Reduction & 2026 & Scopus & kvalificerar eller motsäger påståendet (0.95) \\
Getting the most out of audience response systems: predicting student reactions & 2012 & WoS & kvalificerar eller motsäger påståendet (0.97) \\
How the initial thinking period affects student argumentation during peer instruction: students' experiences versus observations & 2016 & WoS & kvalificerar eller motsäger påståendet (0.96) \\
Impact of Clicker and Confidence Questions on the Metacognition and Performance of Students of Different Achievement Groups in General Chemistry & 2023 & Scopus & kvalificerar eller motsäger påståendet (0.99) \\
Impact of Wireless Technology in a Human Sexuality Course & 2011 & OA & kvalificerar eller motsäger påståendet (0.89) \\
Implementing Audience Response System in Structured Mentoring Processes to Increase Learning Motivation & 2022 & WoS & kvalificerar eller motsäger påståendet (0.82) \\
Incorporation of a game-based approach into the EFL online classrooms: students' perceptions & 2023 & WoS & kvalificerar eller motsäger påståendet (0.96) \\
Integrating Student Response System into English as a Foreign Language Classroom: A Preliminary Study & 2023 & IEEE & kvalificerar eller motsäger påståendet (0.94) \\
Interactive Learning with Student Response System to Encourage Students to Provide Peer Feedback & 2023 & Scopus & kvalificerar eller motsäger påståendet (0.95) \\
Interactivity software tools for teaching in ophthalmology & 2024 & Scopus & kvalificerar eller motsäger påståendet (0.89) \\
INVESTIGATING TECHNIQUE EFFICACY IN EFL READING INSTRUCTION USING PEAR DECK & 2025 & Scopus & kvalificerar eller motsäger påståendet (0.98) \\
Let me explain! The effects of writing and reading short justifications on students' performance, confidence and opinions in audience response systems & 2022 & Scopus & kvalificerar eller motsäger påståendet (0.98) \\
Music, Stories, and Progress Clickers Experiences Improving Classroom Climate with ``Small'' Socio-emotional Activities & 2022 & Scopus & kvalificerar eller motsäger påståendet (0.79) \\
Pre-Service Elementary Teachers' Perceptions of a Modified Peer Instruction Implementation of Clickers in Their Mathematics Content Course. & 2013 & OA & kvalificerar eller motsäger påståendet (0.95) \\
Response switching and self-efficacy in Peer Instruction classrooms & 2015 & OA & kvalificerar eller motsäger påståendet (0.97) \\
Revisiting Clickers: In-Class Questions Followed by At-Home Reflections Are Associated with Higher Student Performance on Related Exam Questions & 2022 & Scopus & kvalificerar eller motsäger påståendet (0.99) \\
Student Experiences with Static versus Adaptive AI Feedback on Self-Explanations in Computer Science Courses & 2026 & Scopus & kvalificerar eller motsäger påståendet (0.96) \\
Student Perceptions towards Using Clickers and Lecture Software Applications in Hospitality Lecture Courses & 2015 & OA & kvalificerar eller motsäger påståendet (0.97) \\
Student Response Systems (SRS): A Game Changer in College Classrooms?-An Explorative Study of Technology-Enhanced Learning & 2026 & Scopus & kvalificerar eller motsäger påståendet (0.90) \\
Student Response Systems in Initial Teacher Education: A Scoping Review of Web-Based Applications & 2023 & Scopus & kvalificerar eller motsäger påståendet (0.82) \\
Students' perceptions on game-based classroom response system in a computer programming course & 2017 & IEEE & kvalificerar eller motsäger påståendet (0.99) \\
Study and Analysis on the Student Response System Adoption: Experimentation in a Programming Course & 2022 & WoS & kvalificerar eller motsäger påståendet (0.97) \\
The Challenges of Using Kahoot! in Teaching and Learning in Higher Education -- A Systematic Review & 2022 & Scopus & kvalificerar eller motsäger påståendet (0.98) \\
The dark side of personal response systems (PRSs): Boredom, feedback, perceived learning, learning satisfaction & 2020 & OA & kvalificerar eller motsäger påståendet (0.99) \\
The Effect of Differing Audience Response System Question Types on Student Attention in the Veterinary Medical Classroom & 2010 & OA & kvalificerar eller motsäger påståendet (0.99) \\
The effects of personal response systems on student engagement and performance on tests & 2014 & OA & kvalificerar eller motsäger påståendet (0.99) \\
The impact of collaborative and individualized student response system strategies on learner motivation, metacognition, and knowledge transfer & 2012 & WoS & kvalificerar eller motsäger påståendet (0.95) \\
The Impact of Personal Response Systems on Students' Learning Performance & 2018 & OA & kvalificerar eller motsäger påståendet (0.91) \\
The Past, Present, and Future of Clickers: A Review & 2024 & WoS & kvalificerar eller motsäger påståendet (0.98) \\
The positive effect of in-class clicker questions on later exams depends on initial student performance level but not question format & 2018 & Scopus & kvalificerar eller motsäger påståendet (0.98) \\
The Psychometric Properties of Classroom Response System Data: A Case Study & 2016 & Scopus & kvalificerar eller motsäger påståendet (0.82) \\
The use of a Classroom Response System to more effectively flip the classroom & 2013 & IEEE & kvalificerar eller motsäger påståendet (0.82) \\
Throtel, throttel, trottel or trotel. Engineering Undergraduates' Perception of the Use of Kahoot! to Review Course Content & 2024 & Scopus & kvalificerar eller motsäger påståendet (0.95) \\
Undergraduate Game-Based Student Response Systems (SRSs): A Systematic Review & 2023 & Scopus & kvalificerar eller motsäger påståendet (0.98) \\
USAGE OF A WEB-BASED STUDENT RESPONSE SYSTEM (SRS) IN THE CLASSROOM: AN ANALYSIS OF ACCOUNTING STUDENTS' PERCEPTION & 2021 & WoS & kvalificerar eller motsäger påståendet (0.97) \\
Use of a classroom response system to enhance classroom interactivity & 2006 & IEEE & kvalificerar eller motsäger påståendet (0.99) \\
Use of Personal Response Systems in Higher Education -- A Case Study & 2014 & OA & kvalificerar eller motsäger påståendet (0.91) \\
Using a personal response system for promoting student interaction & 2003 & IEEE & kvalificerar eller motsäger påståendet (0.82) \\
Using Clickers in Science, Technology, Engineering, and Mathematics & 2023 & Scopus & kvalificerar eller motsäger påståendet (0.91) \\
Using Socrative software for instant formative feedback in physics courses & 2019 & Scopus & kvalificerar eller motsäger påståendet (0.98) \\
Work in progress - `Real World Problems' as assessment of software engineering & 2007 & IEEE & kvalificerar eller motsäger påståendet (0.95) \\
\end{longtable}
\endgroup

